\section{酸塩基の定義と溶媒依存性}

\subsection{平成30年度 問11}
\label{q:h30-reg-applied-11}

\noindent\textbf{原典：} sources/H30/17\_applied\_science.pdf、PDF物理4頁。

\textbf{公式正答：} 4（single）。

\par\textbf{検討未確定：} 4は濃度と温度を指定しない。初回記録のルクス・フラッドに関する短文は電子対と酸化物イオンを混同しており、原記録を保持して本文では定義を訂正する。

\paragraph{問題の要旨} 溶媒系・酸化物イオン・電子対の酸塩基定義と、硫酸の酸化性を区別する。

\Needspace{4\baselineskip}\paragraph{確認した条件・記述（独自要約）}\mbox{}\par

\noindent\textbf{確認範囲：} 問題幹、否定、全5肢、数式、数値、単位を原典画像で確認した。\par

\noindent\textbf{成立条件：} 最も不適切な記述。酸の強さ、酸化還元、溶媒の相対的なプロトン授受を分離する。\par

\paragraph{問われる核心} 溶媒系・酸化物イオン・電子対の酸塩基定義と、硫酸の酸化性を区別する。

\paragraph{必要知識} 最も不適切な記述。酸の強さ、酸化還元、溶媒の相対的なプロトン授受を分離する。

\paragraph{考え方・正答を選ぶ理由（査読待ち初稿）} 溶媒系では自己解離で生じる陽イオンを増やす物質が酸、ルクス・フラッドでは酸化物イオン受容体、ルイスでは電子対受容体が酸である。酢酸もより強い酸である硫酸溶媒中ではプロトンを受ける塩基として働く。硫酸の酸化性は濃度・温度に依存し、一般に非常に強い酸化剤とする4は適切でない。

\Needspace{4\baselineskip}\paragraph{選択肢の検討（査読待ち初稿）}\mbox{}\par

\noindent\textbf{1：適切} 溶媒の自己解離陽イオンを生じさせる物質は酸。溶媒系の定義であり、水に限らない。\par

\noindent\textbf{2：適切} 酸化物イオンを受け取る物質は酸。ルクス・フラッドの定義である。電子対受容の定義とは区別する。\par

\noindent\textbf{3：適切} 電子対を受け取る物質は酸。ルイスの定義である。\par

\noindent\textbf{4：不適切} 硫酸は非常に強い酸化剤。希硫酸と熱濃硫酸を同じ酸化剤として一括できない。酸性の強さから酸化性を導けない。\par

\noindent\textbf{5：適切} 酢酸は水中で酸、硫酸中で塩基。強酸溶媒中のプロトン受容という相対的役割である。\par

\paragraph{再利用する考え方} 反応速度と平衡位置、全圧と分圧、濃度と活量を分ける。式の適用条件を確認し、原典の不正確な記述を科学的に正しい文へ置き換えない。

\paragraph{共通理論} \hyperref[theory:chem_eq:definitions]{第\ref*{theory:chem_eq:definitions}節}

\par 参考：\cite{iupac_greenbook2012}。

\paragraph{関連問題} \hyperref[q:r06-reg-applied-12]{R06-REG-APPLIED-12}

\par 難易度 2（校準初稿）、優先度未評価。

\subsection{令和6年度 問12}
\label{q:r06-reg-applied-12}

\noindent\textbf{原典：} sources/R06/17\_applied\_science.pdf、PDF物理4頁。

\textbf{公式正答：} 5（single）。

\par\textbf{検討未確定：} 3の原典表現は酸解離定数pKaである。pKaとKaを同一視する文を科学的に修正した原文として保存せず、用語の不正確さを注記する。

\paragraph{問題の要旨} ルイス・ブレンステッド・アレニウスの酸、pKa、塩素のオキソ酸の強さを比較する。

\Needspace{4\baselineskip}\paragraph{確認した条件・記述（独自要約）}\mbox{}\par

\noindent\textbf{確認範囲：} 問題幹、否定、全5肢、数式、数値、単位を原典画像で確認した。\par

\noindent\textbf{成立条件：} 酸解離定数は活量に基づく無次元量でありpKaはその負常用対数。HClOとHClO2は同じ中心元素で比較。\par

\paragraph{問われる核心} ルイス・ブレンステッド・アレニウスの酸、pKa、塩素のオキソ酸の強さを比較する。

\paragraph{必要知識} 酸解離定数は活量に基づく無次元量でありpKaはその負常用対数。HClOとHClO2は同じ中心元素で比較。

\paragraph{考え方・正答を選ぶ理由（査読待ち初稿）} ルイス酸は電子対受容体、ブレンステッド酸はプロトン供与体、アレニウス酸は水中で水素イオンを増やす物質である。\(\mathrm{p}K_a=-\log_{10}K_a\) なので小さいほど強い酸となる。塩素中心のオキソ酸では亜塩素酸HClO2が次亜塩素酸HClOより強く、逆を述べる5が誤りである。

\Needspace{4\baselineskip}\paragraph{選択肢の検討（査読待ち初稿）}\mbox{}\par

\noindent\textbf{1：適切} ルイス酸は電子対受容体。共有結合形成に電子対を受け取る。\par

\noindent\textbf{2：適切} ブレンステッド酸はプロトン供与体。相手の塩基にプロトンを渡す。\par

\noindent\textbf{3：条件付きで適切} 酸解離定数pKaが小さいほど強酸。数値の大小関係はpKaについて正しい。pKaそのものを酸解離定数と呼ぶ点は不正確である。\par

\noindent\textbf{4：適切} アレニウス酸は水中で水素イオンを生じる。水溶液に限定した定義である。\par

\noindent\textbf{5：不適切} 次亜塩素酸は亜塩素酸より強い酸。Purdueの表ではHOClのKaが2.9の10のマイナス8乗、HOClOが1.1の10のマイナス2乗であり、逆である。\par

\paragraph{再利用する考え方} 反応速度と平衡位置、全圧と分圧、濃度と活量を分ける。式の適用条件を確認し、原典の不正確な記述を科学的に正しい文へ置き換えない。

\paragraph{共通理論} \hyperref[theory:chem_eq:definitions]{第\ref*{theory:chem_eq:definitions}節}, \hyperref[theory:chem_eq:acid_strength]{第\ref*{theory:chem_eq:acid_strength}節}

\par 参考：\cite{iupac_greenbook2012,purdue_acid_strength}。

\paragraph{関連問題} \hyperref[q:h30-reg-applied-11]{H30-REG-APPLIED-11}

\par 難易度 2（校準初稿）、優先度未評価。

\section{分離操作の原理と共沸による蒸留限界を判断する。}

\subsection{平成30年度 問18}
\label{q:h30-reg-applied-18}

\noindent\textbf{原典：} sources/H30/17\_applied\_science.pdf、PDF物理6頁。

\textbf{公式正答：} 1（single）。

\par\textbf{検討未確定：} 逆浸透膜の全分子・イオンの完全阻止という記述は理想化。公式は蒸留の共沸例外を選ぶが、実膜の阻止率と対象化学種の条件を留保。

\paragraph{問題の要旨} 分離操作の原理と共沸による蒸留限界を判断する。

\Needspace{4\baselineskip}\paragraph{確認した条件・記述（独自要約）}\mbox{}\par

\noindent\textbf{原典確認：} S7\_049の26原典画像で全五肢・否定・添字を確認。\par

\noindent\textbf{対象：} 分離操作の原理と共沸による蒸留限界を判断する。\par

\paragraph{問われる核心} 分離操作の原理と共沸による蒸留限界を判断する。

\paragraph{必要知識} 蒸留は気液平衡の組成差を使う。共沸点でy=\allowbreak{}xとなる例があるため、必ず異なるという全称が誤りである。吸収は液体内部への溶解、吸着は表面や細孔への保持、抽出は非混和液相間の分配として区別する。逆浸透は圧力で水を透過させる代表的操作だが、理想的な完全溶質阻止をすべての膜に保証しない。

\paragraph{考え方・正答を選ぶ理由（査読待ち初稿）} 蒸留は気液平衡の組成差を使う。共沸点でy=\allowbreak{}xとなる例があるため、必ず異なるという全称が誤りである。吸収は液体内部への溶解、吸着は表面や細孔への保持、抽出は非混和液相間の分配として区別する。逆浸透は圧力で水を透過させる代表的操作だが、理想的な完全溶質阻止をすべての膜に保証しない。

\Needspace{4\baselineskip}\paragraph{選択肢の検討（査読待ち初稿）}\mbox{}\par

\noindent\textbf{1：不適切} 蒸留で気液組成は必ず違う。共沸組成では気液の組成が一致し、通常の単純な蒸留を反復してもその組成を越えて分離できない。\par

\noindent\textbf{2：通常の条件で適切} 気体成分の液体への吸収。気体を液体と接触させ、溶けやすい成分を液相へ移す。表面への吸着と区別する。\par

\noindent\textbf{3：通常の条件で適切} 非混和液相への抽出。二相間の分配平衡を利用し、溶質を片方の相へ移す。\par

\noindent\textbf{4：通常の条件で適切} 活性炭・シリカゲル・ゼオライト。表面や細孔への相互作用により成分を保持する吸着剤として用いられる。\par

\noindent\textbf{5：条件付き・解釈留保} 逆浸透膜の水の選択透過。溶媒を選択的に透過させ、浸透圧差に逆らう圧力で濃縮・脱塩する。実膜は全溶質を100\%阻止する理想膜とは限らない。\par

\paragraph{再利用する考え方} 測定対象と分離対象、定義の分母、絶対量と相対校正、全称と成立条件を分けて判断する。

\paragraph{共通理論} \hyperref[theory:analysis:operations]{第\ref*{theory:analysis:operations}節}

\par 参考：\cite{nptel_azeotrope,inorganic_zeolite}。

\paragraph{関連問題} \hyperref[q:r02-reg-applied-17]{R02-REG-APPLIED-17}

\par 難易度 4（校準初稿）、優先度未評価。

\subsection{令和2年度 問17}
\label{q:r02-reg-applied-17}

\noindent\textbf{原典：} sources/R02/17\_applied\_science.pdf、PDF物理6頁。

\textbf{公式正答：} 4（single）。

\par\textbf{検討未確定：} 逆浸透膜の全分子・イオンの完全阻止という記述は理想化。公式は蒸留の共沸例外を選ぶが、実膜の阻止率と対象化学種の条件を留保。

\paragraph{問題の要旨} 分離操作の原理と共沸による蒸留限界を判断する。

\Needspace{4\baselineskip}\paragraph{確認した条件・記述（独自要約）}\mbox{}\par

\noindent\textbf{原典確認：} S7\_049の26原典画像で全五肢・否定・添字を確認。\par

\noindent\textbf{対象：} 分離操作の原理と共沸による蒸留限界を判断する。\par

\paragraph{問われる核心} 分離操作の原理と共沸による蒸留限界を判断する。

\paragraph{必要知識} 蒸留は気液平衡の組成差を使う。共沸点でy=\allowbreak{}xとなる例があるため、必ず異なるという全称が誤りである。吸収は液体内部への溶解、吸着は表面や細孔への保持、抽出は非混和液相間の分配として区別する。逆浸透は圧力で水を透過させる代表的操作だが、理想的な完全溶質阻止をすべての膜に保証しない。

\paragraph{考え方・正答を選ぶ理由（査読待ち初稿）} 蒸留は気液平衡の組成差を使う。共沸点でy=\allowbreak{}xとなる例があるため、必ず異なるという全称が誤りである。吸収は液体内部への溶解、吸着は表面や細孔への保持、抽出は非混和液相間の分配として区別する。逆浸透は圧力で水を透過させる代表的操作だが、理想的な完全溶質阻止をすべての膜に保証しない。

\Needspace{4\baselineskip}\paragraph{選択肢の検討（査読待ち初稿）}\mbox{}\par

\noindent\textbf{1：通常の条件で適切} 活性炭・シリカゲル・ゼオライト。表面や細孔への相互作用により成分を保持する吸着剤として用いられる。\par

\noindent\textbf{2：通常の条件で適切} 非混和液相への抽出。二相間の分配平衡を利用し、溶質を片方の相へ移す。\par

\noindent\textbf{3：通常の条件で適切} 気体成分の液体への吸収。気体を液体と接触させ、溶けやすい成分を液相へ移す。表面への吸着と区別する。\par

\noindent\textbf{4：不適切} 蒸留で気液組成は必ず違う。共沸組成では気液の組成が一致し、通常の単純な蒸留を反復してもその組成を越えて分離できない。\par

\noindent\textbf{5：条件付き・解釈留保} 逆浸透膜の水の選択透過。溶媒を選択的に透過させ、浸透圧差に逆らう圧力で濃縮・脱塩する。実膜は全溶質を100\%阻止する理想膜とは限らない。\par

\paragraph{再利用する考え方} 測定対象と分離対象、定義の分母、絶対量と相対校正、全称と成立条件を分けて判断する。

\paragraph{共通理論} \hyperref[theory:analysis:operations]{第\ref*{theory:analysis:operations}節}

\par 参考：\cite{nptel_azeotrope,inorganic_zeolite}。

\paragraph{関連問題} \hyperref[q:h30-reg-applied-18]{H30-REG-APPLIED-18}

\par 難易度 4（校準初稿）、優先度未評価。

\section{溶液分母の溶解度と析出収支}

\subsection{平成27年度 問11}
\label{q:h27-reg-applied-11}

\noindent\textbf{原典：} sources/H27/17\_applied\_science.pdf、PDF物理4頁。

\textbf{公式正答：} 4（single）。

\par\textbf{検討未確定：} 水の蒸発なし・析出結晶への水の取込みなしを置く物質収支。原典の分母は飽和溶液。

\paragraph{問題の要旨} 飽和溶液質量を分母にした溶解度から析出量を求める。

\Needspace{4\baselineskip}\paragraph{確認した条件・記述（独自要約）}\mbox{}\par

\noindent\textbf{60℃溶質：} 飽和溶液100g中52.2g\par

\noindent\textbf{20℃溶質：} 飽和溶液100g中24.0g\par

\noindent\textbf{水の初期量：} 47.8g\par

\paragraph{問われる核心} 飽和溶液質量を分母にした溶解度から析出量を求める。

\paragraph{必要知識} 原典の定義は100gの飽和水溶液中の溶質量。

\paragraph{考え方・正答を選ぶ理由（査読待ち初稿）} 原典の定義は100gの飽和水溶液中の溶質量。60℃の100gには溶質52.2g、水47.8gがある。蒸発を無視し20℃で溶質/水=24/76とすると残る溶質は47.8×24/76=\allowbreak{}15.0947g。析出52.2−15.0947=\allowbreak{}37.1053gで4。100g水当たりという通常の別定義へ読み替えない。

\Needspace{4\baselineskip}\paragraph{選択肢の検討（査読待ち初稿）}\mbox{}\par

\noindent\textbf{1：選ばない候補} 15.1g。冷却後の溶存量で析出量ではない。\par

\noindent\textbf{2：選ばない候補} 24.0g。冷却前の水量を100gと誤認。\par

\noindent\textbf{3：選ばない候補} 28.2g。52.2−24.0は溶液量が変わることを無視。\par

\noindent\textbf{4：選ぶ候補} 37.1g。水47.8gを保持し溶存15.0947gを差し引く。\par

\noindent\textbf{5：選ばない候補} 59.0g。初期溶質52.2gより大きく質量保存に反する。\par

\paragraph{再利用する考え方} 軌道と電子状態、格子と基底、単位胞原子数と配位数、定義の分母を区別して候補を全て評価する。

\paragraph{共通理論} \hyperref[theory:atomic:solubility]{第\ref*{theory:atomic:solubility}節}

\par 難易度 3（校準初稿）、優先度未評価。

\section{ボルンハーバーサイクル}

\subsection{平成26年度 問15}
\label{q:h26-reg-applied-15}

\noindent\textbf{原典：} sources/H26/17\_applied\_science.pdf、PDF物理6頁。

\textbf{公式正答：} 2（single）。

\par\textbf{検討未確定：} 原典は「生成熱」を負の候補で表す一方、他の熱量は正の大きさで与える。現代的な生成エンタルピーと熱化学式の右辺の発熱量の符号を区別して説明する。

\paragraph{問題の要旨} ボルン・ハーバーサイクルの各過程に符号を付けてNaClの生成エンタルピーを求める。

\Needspace{4\baselineskip}\paragraph{確認した条件・記述（独自要約）}\mbox{}\par

\noindent\textbf{確認範囲：} 問題幹、数値、式、否定、全5候補を原典画像で照合した。\par

\noindent\textbf{成立条件：} 原典の電子親和力354と格子エネルギー788は正の大きさとして与えられている。電子獲得と結晶生成はエンタルピー低下であり負符号、Cl原子1 molにはCl2を半mol解離する。\par

\paragraph{問われる核心} ボルン・ハーバーサイクルの各過程に符号を付けてNaClの生成エンタルピーを求める。

\paragraph{必要知識} 原典の電子親和力354と格子エネルギー788は正の大きさとして与えられている。電子獲得と結晶生成はエンタルピー低下であり負符号、Cl原子1 molにはCl2を半mol解離する。

\paragraph{考え方・正答を選ぶ理由（査読待ち初稿）} Na(s)の昇華+108、Na(g)のイオン化+502、Cl2(g)半molの解離+243/2、Cl(g)の電子獲得−354、気体イオンからNaCl(s)の生成−788を加える。\(108+502+121.5-354-788=-410.5\,\mathrm{kJ\,mol^{-1}}\)。最も近い−411 kJ/molの2を選ぶ。

\Needspace{4\baselineskip}\paragraph{選択肢の検討（査読待ち初稿）}\mbox{}\par

\noindent\textbf{1：不適切} −289 kJ/mol。Cl2の解離熱243を半分にせず全量使うと−289になる。NaCl 1 molに必要なCl原子は1 molである。\par

\noindent\textbf{2：適切} −411 kJ/mol。半molの塩素分子解離と吸熱・発熱の方向をすべて反映する。\par

\noindent\textbf{3：不適切} −471 kJ/mol。指定された5過程のエンタルピー和−410.5と合わない。\par

\noindent\textbf{4：不適切} −597 kJ/mol。原典の昇華、イオン化、半解離、電子獲得、格子形成の和と合わない。\par

\noindent\textbf{5：不適切} −1166 kJ/mol。吸熱の昇華・イオン化・解離を含めた生成エンタルピーより過剰な負値である。\par

\paragraph{再利用する考え方} 電子親和力と格子エネルギーは文献で符号慣習が異なる。名称だけに符号を割り当てず、反応方向を書いてから足す。

\paragraph{共通理論} \hyperref[theory:electrochem:born_haber]{第\ref*{theory:electrochem:born_haber}節}, \hyperref[theory:electrochem:heat_sign]{第\ref*{theory:electrochem:heat_sign}節}

\par 参考：\cite{openstax_born_haber}。

\paragraph{関連問題} \hyperref[q:h24-reg-applied-14]{H24-REG-APPLIED-14}

\par 難易度 2（校準初稿）、優先度未評価。

\section{共役酸塩基対・緩衝容量・pH}

\subsection{令和3年度 問17}
\label{q:r03-reg-applied-17}

\noindent\textbf{原典：} sources/R03/17\_applied\_science.pdf、PDF物理6頁。

\textbf{公式正答：} 2（single）。

\par\textbf{検討未確定：} 5は混合比を指定しない。部分中和という成立条件を注記し、一律に全混合比で適切とは扱わない。

\paragraph{問題の要旨} 緩衝溶液の定義、血液、式、酸塩基対の調製条件を判定する。

\Needspace{4\baselineskip}\paragraph{確認した条件・記述（独自要約）}\mbox{}\par

\noindent\textbf{確認範囲：} 問題幹、否定、全5肢、数式、数値、単位を原典画像で確認した。\par

\noindent\textbf{成立条件：} 最も不適切な文。緩衝は有限量の酸塩基添加に対する変化の小ささであり、pH完全不変ではない。弱塩基の強酸による部分中和。\par

\paragraph{問われる核心} 緩衝溶液の定義、血液、式、酸塩基対の調製条件を判定する。

\paragraph{必要知識} 最も不適切な文。緩衝は有限量の酸塩基添加に対する変化の小ささであり、pH完全不変ではない。弱塩基の強酸による部分中和。

\paragraph{考え方・正答を選ぶ理由（査読待ち初稿）} 緩衝溶液には弱酸と共役塩基、又は弱塩基と共役酸の両方が必要である。血液のpH維持にも緩衝作用が関与するので、それを否定する2が誤りである。\(K_a=a_{\mathrm{H^+}}a_{\mathrm{A^-}}/a_{\mathrm{HA}}\) を対数化すると \(\mathrm{pH=p}K_a+\log_{10}(a_{\mathrm{A^-}}/a_{\mathrm{HA}})\) となる。弱塩基と強酸の混合は部分中和で両成分を残す条件が必要である。

\Needspace{4\baselineskip}\paragraph{選択肢の検討（査読待ち初稿）}\mbox{}\par

\noindent\textbf{1：適切} 少量の強酸強塩基添加でpHが大きく変わらない。緩衝容量内での性質である。\par

\noindent\textbf{2：不適切} 血液のpH7.4維持に緩衝は関与しない。重炭酸系やタンパク等の酸塩基対が関与する。\par

\noindent\textbf{3：条件付きで適切} ヘンダーソン・ハッセルバルヒ式でpHを求める。活量比の式は平衡から導かれる。濃度比や混合時の物質量比に置換するには近似条件が必要である。\par

\noindent\textbf{4：適切} 弱酸とその共役塩基の塩で調製。両成分が有効量存在するように調製する。\par

\noindent\textbf{5：条件付きで適切} 弱塩基と強酸で緩衝溶液を作る。部分中和で弱塩基と共役酸の両方を残す。強酸過剰では一般に緩衝溶液にならない。\par

\paragraph{再利用する考え方} 反応速度と平衡位置、全圧と分圧、濃度と活量を分ける。式の適用条件を確認し、原典の不正確な記述を科学的に正しい文へ置き換えない。

\paragraph{共通理論} \hyperref[theory:chem_eq:buffer]{第\ref*{theory:chem_eq:buffer}節}

\par 参考：\cite{iupac_greenbook2012,mit_buffer2008}。

\par 難易度 2（校準初稿）、優先度未評価。

\section{触媒・反応経路・平衡不変}

\subsection{平成23年度 問12}
\label{q:h23-reg-applied-12}

\noindent\textbf{原典：} sources/H23/17\_applied\_science.pdf、PDF物理5頁。

\textbf{公式正答：} 3（single）。

\par\textbf{検討未確定：} 負触媒という出題時の用語と現行IUPAC定義の不一致、及び自ら変化せずの正味再生という省略を保持する。

\paragraph{問題の要旨} 触媒が変える反応速度と、変えない平衡定数・反応熱を区別する。

\Needspace{4\baselineskip}\paragraph{確認した条件・記述（独自要約）}\mbox{}\par

\noindent\textbf{確認範囲：} 問題幹、否定、全5肢、数式、数値、単位を原典画像で確認した。\par

\noindent\textbf{成立条件：} 触媒の正味の再生、同一温度での平衡不変、反応熱の状態関数性。\par

\paragraph{問われる核心} 触媒が変える反応速度と、変えない平衡定数・反応熱を区別する。

\paragraph{必要知識} 触媒の正味の再生、同一温度での平衡不変、反応熱の状態関数性。

\paragraph{考え方・正答を選ぶ理由（査読待ち初稿）} 触媒は別の反応経路を与えるが、反応前後の状態が同じなら反応熱を変えない。したがって反応熱の減少を述べる3が不適切である。4の負触媒は歴史的用語として読み、現行IUPACでは速度を下げる作用を阻害として区別する。

\Needspace{4\baselineskip}\paragraph{選択肢の検討（査読待ち初稿）}\mbox{}\par

\noindent\textbf{1：条件付きで適切} 自ら変化せず速度を変え又は反応を開始する。反応過程では中間体として変化し得るが、正味で再生されるという意味である。開始は実用上観測できる速度への増大として読む。\par

\noindent\textbf{2：適切} 平衡に影響しない。同じ温度で反応の標準ギブズエネルギー差を変えない。\par

\noindent\textbf{3：不適切} 反応熱を減少させる。活性化障壁の低下と反応前後のエンタルピー差は別である。\par

\noindent\textbf{4：条件付きで適切} 負触媒は速度を低下させる。歴史的呼称ではその意味であるが、IUPACは負触媒という用語を採用せず阻害を区別する。\par

\noindent\textbf{5：適切} 選択性を持ち特定生成物を与える。競合する反応経路の相対速度を変えて生成物選択性を持ち得る。\par

\paragraph{再利用する考え方} 反応速度と平衡位置、全圧と分圧、濃度と活量を分ける。式の適用条件を確認し、原典の不正確な記述を科学的に正しい文へ置き換えない。

\paragraph{共通理論} \hyperref[theory:chem_eq:catalyst]{第\ref*{theory:chem_eq:catalyst}節}

\par 参考：\cite{iupac_greenbook2012,iupac_catalyst}。

\paragraph{関連問題} \hyperref[q:r01-ret-applied-15]{R01-RET-APPLIED-15}

\par 難易度 2（校準初稿）、優先度未評価。

\subsection{令和元年度再試験 問15}
\label{q:r01-ret-applied-15}

\noindent\textbf{原典：} sources/R01\_retest/17\_applied\_science.pdf、PDF物理6頁。

\textbf{公式正答：} 5（single）。

\paragraph{問題の要旨} 活性化障壁・動的平衡・触媒選択性と、速度と平衡位置の独立性を問う。

\Needspace{4\baselineskip}\paragraph{確認した条件・記述（独自要約）}\mbox{}\par

\noindent\textbf{確認範囲：} 問題幹、否定、全5肢、数式、数値、単位を原典画像で確認した。\par

\noindent\textbf{成立条件：} 温度・圧力一定。同じ反応物と生成物について平衡と速度を比較する。活性化エネルギーは問題の反応経路上の障壁模型として読む。\par

\paragraph{問われる核心} 活性化障壁・動的平衡・触媒選択性と、速度と平衡位置の独立性を問う。

\paragraph{必要知識} 温度・圧力一定。同じ反応物と生成物について平衡と速度を比較する。活性化エネルギーは問題の反応経路上の障壁模型として読む。

\paragraph{考え方・正答を選ぶ理由（査読待ち初稿）} 触媒は障壁の低い別経路を与えるが、温度一定で平衡定数を変えない。平衡では正逆の反応が停止するのではなく同じ速度で継続する。定温定圧で反応速度が増すと平衡位置も変わるとする5が誤りである。

\Needspace{4\baselineskip}\paragraph{選択肢の検討（査読待ち初稿）}\mbox{}\par

\noindent\textbf{1：適切} 反応物から遷移状態への差が障壁。指定した反応経路のエネルギー障壁を説明する。経験的アレニウス量と無条件に同一視しない。\par

\noindent\textbf{2：適切} 平衡では正逆反応速度が等しい。濃度の正味変化が0でも各方向の反応は進む。\par

\noindent\textbf{3：適切} 触媒は低い障壁の経路を与える。同じ反応の別経路を通じて速度を上げる。\par

\noindent\textbf{4：適切} 触媒は選択性を持つ。競合経路の相対速度を変え得る。\par

\noindent\textbf{5：不適切} 速度増大で定温定圧の平衡位置が変わる。速度は平衡への到達時間に関わり、平衡位置は自由エネルギーで決まる。\par

\paragraph{再利用する考え方} 反応速度と平衡位置、全圧と分圧、濃度と活量を分ける。式の適用条件を確認し、原典の不正確な記述を科学的に正しい文へ置き換えない。

\paragraph{共通理論} \hyperref[theory:chem_eq:catalyst]{第\ref*{theory:chem_eq:catalyst}節}, \hyperref[theory:chem_eq:activation]{第\ref*{theory:chem_eq:activation}節}, \hyperref[theory:chem_eq:rates]{第\ref*{theory:chem_eq:rates}節}

\par 参考：\cite{iupac_greenbook2012}。

\paragraph{関連問題} \hyperref[q:h23-reg-applied-12]{H23-REG-APPLIED-12}

\par 難易度 2（校準初稿）、優先度未評価。

\section{束一的性質・凝固点降下}

\subsection{平成24年度 問11}
\label{q:h24-reg-applied-11}

\noindent\textbf{原典：} sources/H24/17\_applied\_science.pdf、PDF物理3頁。

\textbf{公式正答：} 4（single）。

\paragraph{問題の要旨} 完全解離する食塩の溶液について、粒子数を数えて凝固点降下を計算する。

\Needspace{4\baselineskip}\paragraph{確認した条件・記述（独自要約）}\mbox{}\par

\noindent\textbf{確認範囲：} 問題幹、条件、否定、数値、単位、全5肢を原典画像で照合した。\par

\noindent\textbf{成立条件：} 水1.0 kg、NaCl 0.20 mol、\(K_f=1.85\,\mathrm{K\,kg\,mol^{-1}}\)、完全イオン解離の指定。希薄溶液の束一的性質の式を用いる。\par

\paragraph{問われる核心} 完全解離する食塩の溶液について、粒子数を数えて凝固点降下を計算する。

\paragraph{必要知識} 水1.0 kg、NaCl 0.20 mol、\(K_f=1.85\,\mathrm{K\,kg\,mol^{-1}}\)、完全イオン解離の指定。希薄溶液の束一的性質の式を用いる。

\paragraph{考え方・正答を選ぶ理由（査読待ち初稿）} 溶媒の水は1.0 kg、NaClは0.20 molである。完全解離によりNaとClのイオンが各0.20 molとなり、粒子基準の質量モル濃度は \(2\times0.20/1.0=0.40\,\mathrm{mol\,kg^{-1}}\) である。正の凝固点降下量は \(\Delta T_f=1.85\times0.40=0.74\,\mathrm K\)。純水の凝固点を0°Cとして溶液は−0.74°Cとなり4を選ぶ。

\Needspace{4\baselineskip}\paragraph{選択肢の検討（査読待ち初稿）}\mbox{}\par

\noindent\textbf{1：不適切} −0.19°C。完全解離後の粒子濃度0.40 mol/kgを使った降下量0.74 Kと合わない。\par

\noindent\textbf{2：不適切} −0.37°C。NaClの式量単位だけを数え、2種類のイオンへの解離を無視した値である。\par

\noindent\textbf{3：不適切} −0.56°C。粒子濃度と係数の積0.74 Kに合わない。\par

\noindent\textbf{4：適切} −0.74°C。解離による粒子数2倍と凝固点の低下方向を共に反映する。\par

\noindent\textbf{5：不適切} −1.9°C。指定濃度は粒子基準でも0.40 mol/kgであり、1 mol/kg相当の降下とは異なる。\par

\paragraph{再利用する考え方} 質量モル濃度の分母は溶液体積でなく溶媒質量である。完全解離の仮定と、濃厚電解質の実際の活量補正を区別する。

\paragraph{共通理論} \hyperref[theory:phase_sep:freezing]{第\ref*{theory:phase_sep:freezing}節}

\par 参考：\cite{mit_colligative2008}。

\par 難易度 2（校準初稿）、優先度未評価。

\section{衝突・反応断面積・拡散律速}

\subsection{令和元年度 問15}
\label{q:r01-reg-applied-15}

\noindent\textbf{原典：} sources/R01/17\_applied\_science.pdf、PDF物理5頁。

\textbf{公式正答：} 3（single）。

\par\textbf{検討未確定：} 衝突頻度の対象、立体因子模型、アレニウス活性化エネルギーを障壁差で代表する省略を原典通り保持し、適用範囲を教材で限定する。

\paragraph{問題の要旨} 衝突模型、拡散律速、平衡の熱力学、活性化エネルギーの意味を比較する。

\Needspace{4\baselineskip}\paragraph{確認した条件・記述（独自要約）}\mbox{}\par

\noindent\textbf{確認範囲：} 問題幹、否定、全5肢、数式、数値、単位を原典画像で確認した。\par

\noindent\textbf{成立条件：} 気体圧力比較は温度・組成一定。衝突頻度の単位と対象を明記する。アレニウス式の経験的活性化エネルギーと電子的障壁高さは一般には別。\par

\paragraph{問われる核心} 衝突模型、拡散律速、平衡の熱力学、活性化エネルギーの意味を比較する。

\paragraph{必要知識} 気体圧力比較は温度・組成一定。衝突頻度の単位と対象を明記する。アレニウス式の経験的活性化エネルギーと電子的障壁高さは一般には別。

\paragraph{考え方・正答を選ぶ理由（査読待ち初稿）} 出会った後のプロトン移動が十分速い中和では、分子やイオンが出会う拡散が速度を制限する。反応そのものが律速とする3が不適切である。1は温度と組成一定の分子1個あたりの衝突頻度、2は単純な立体因子模型、5は障壁模型という適用範囲を持つ。

\Needspace{4\baselineskip}\paragraph{選択肢の検討（査読待ち初稿）}\mbox{}\par

\noindent\textbf{1：条件付きで適切} 衝突頻度は圧力に比例。分子1個あたりの頻度は数密度に比例する。単位体積あたりの衝突数なら固定組成では圧力の2乗に比例する。\par

\noindent\textbf{2：条件付きで適切} 反応断面積は衝突断面積と立体因子の積。配向の割合を掛ける単純模型である。一般の反応断面積の定義を尽くさない。\par

\noindent\textbf{3：不適切} 速い酸塩基中和は反応律速。遭遇後の反応が速い極限では拡散による遭遇が律速する。\par

\noindent\textbf{4：適切} 標準自由エネルギー変化は平衡定数を決める。\(K=\exp(-\Delta_rG^\circ/RT)\) に従う。\par

\noindent\textbf{5：条件付きで適切} アレニウス活性化エネルギーは活性錯合体との差。単純な障壁模型の説明である。厳密な経験的定義は \(E_a=RT^2d\ln k/dT\) であり、常に障壁高さと同じではない。\par

\paragraph{再利用する考え方} 反応速度と平衡位置、全圧と分圧、濃度と活量を分ける。式の適用条件を確認し、原典の不正確な記述を科学的に正しい文へ置き換えない。

\paragraph{共通理論} \hyperref[theory:chem_eq:collision]{第\ref*{theory:chem_eq:collision}節}, \hyperref[theory:chem_eq:diffusion]{第\ref*{theory:chem_eq:diffusion}節}, \hyperref[theory:chem_eq:activation]{第\ref*{theory:chem_eq:activation}節}

\par 参考：\cite{iupac_greenbook2012}。

\paragraph{関連問題} \hyperref[q:r02-reg-applied-13]{R02-REG-APPLIED-13}

\par 難易度 2（校準初稿）、優先度未評価。

\subsection{令和2年度 問13}
\label{q:r02-reg-applied-13}

\noindent\textbf{原典：} sources/R02/17\_applied\_science.pdf、PDF物理5頁。

\textbf{公式正答：} 2（single）。

\par\textbf{独立解答と公式正答に相違あり：} 独立解答 3。初回3、公式2の不一致を保持する。原典3の速度比は一般に誤りであり、2も比較する断面積の定義に依存する。最も不適切を一意に説明する根拠は未解決であり、公式に合わせて3を速度定数に書き換えない。

\paragraph{問題の要旨} 衝突断面積と反応断面積、反応速度と速度定数、拡散律速の定義を区別する。

\Needspace{4\baselineskip}\paragraph{確認した条件・記述（独自要約）}\mbox{}\par

\noindent\textbf{確認範囲：} 問題幹、否定、全5肢、数式、数値、単位を原典画像で確認した。\par

\noindent\textbf{成立条件：} 原典3は逆反応の反応速度と正反応の反応速度の比。速度定数へ置換しない。断面積は入射流束に対する有効面積であり、比較する衝突の定義とエネルギーを指定する。\par

\paragraph{問われる核心} 衝突断面積と反応断面積、反応速度と速度定数、拡散律速の定義を区別する。

\paragraph{必要知識} 原典3は逆反応の反応速度と正反応の反応速度の比。速度定数へ置換しない。断面積は入射流束に対する有効面積であり、比較する衝突の定義とエネルギーを指定する。

\paragraph{考え方・正答を選ぶ理由（査読待ち初稿）} 1は分子1個あたりの \(z=\sigma\bar v_{\mathrm{rel}}n\)、4は活性化模型、5は拡散律速を述べる。3の原文は速度定数ではなく反応速度の比である。平衡では正逆の速度が等しく比は1なので、一般の平衡定数を表せず、初回は3を選んだ。公式正答は2である。2について剛体球の幾何学断面積と実効的反応断面積を区別する必要があるが、断面積の比較条件は原典に明示されない。2の公式判定だけで3の記述を適切に直すことはできない。

\Needspace{4\baselineskip}\paragraph{選択肢の検討（査読待ち初稿）}\mbox{}\par

\noindent\textbf{1：適切} 衝突頻度は断面積・相対速さ・数密度に比例。分子1個あたりの衝突頻度の式である。\par

\noindent\textbf{2：解釈未解決} 反応断面積は衝突断面積を上回らない。同じ遭遇集合のうち反応する割合という模型なら成立する。剛体球幾何学面積との比較を一般則とすることはできない。公式2の根拠となる断面積の意味は追加査読が必要である。\par

\noindent\textbf{3：不適切} 平衡定数は逆反応速度と正反応速度の比。動的平衡では速度比は1である。一次可逆反応の速度定数比でも順序は正方向の定数を逆方向で割る。\par

\noindent\textbf{4：条件付きで適切} 二分子速度定数は障壁を超える衝突割合に依存。活性化障壁を持つ単純衝突模型である。無障壁反応や拡散律速を包括する一般式ではない。\par

\noindent\textbf{5：適切} 遭遇対が分離より速く生成物を作るなら拡散で決まる。出会った後の反応確率が高い拡散律速の極限である。\par

\paragraph{再利用する考え方} 反応速度と平衡位置、全圧と分圧、濃度と活量を分ける。式の適用条件を確認し、原典の不正確な記述を科学的に正しい文へ置き換えない。

\paragraph{共通理論} \hyperref[theory:chem_eq:collision]{第\ref*{theory:chem_eq:collision}節}, \hyperref[theory:chem_eq:diffusion]{第\ref*{theory:chem_eq:diffusion}節}, \hyperref[theory:chem_eq:rates]{第\ref*{theory:chem_eq:rates}節}

\par 参考：\cite{iupac_greenbook2012}。

\paragraph{関連問題} \hyperref[q:r01-reg-applied-15]{R01-REG-APPLIED-15}

\par 難易度 2（校準初稿）、優先度未評価。

\section{コロイド分散・会合・乳化}

\subsection{令和6年度 問11}
\label{q:r06-reg-applied-11}

\noindent\textbf{原典：} sources/R06/17\_applied\_science.pdf、PDF物理4頁。

\textbf{公式正答：} 2（single）。

\par\textbf{検討未確定：} 溶解という語の範囲、石けんの濃度・温度・ミセル形成条件を原典は限定しない。公式2と初回2を保持し、全称の説明はしない。

\paragraph{問題の要旨} コロイドの分散と分子・イオンの溶解を、乳濁液・光散乱・石けんで区別する。

\Needspace{4\baselineskip}\paragraph{確認した条件・記述（独自要約）}\mbox{}\par

\noindent\textbf{確認範囲：} 問題幹、条件、否定、数値、単位、全5肢を原典画像で照合した。\par

\noindent\textbf{成立条件：} 最も関係が薄い用語の選択。石けんの会合は濃度・温度に依存し、通常の分子・イオンの溶液とコロイド分散を区別する。\par

\paragraph{問われる核心} コロイドの分散と分子・イオンの溶解を、乳濁液・光散乱・石けんで区別する。

\paragraph{必要知識} 最も関係が薄い用語の選択。石けんの会合は濃度・温度に依存し、通常の分子・イオンの溶液とコロイド分散を区別する。

\paragraph{考え方・正答を選ぶ理由（査読待ち初稿）} エマルションは液滴の分散、チンダル現象は分散粒子による光散乱、マヨネーズは乳濁液の具体例、石けん水はミセル等の会合体を含み得る。通常の分子・イオンとして均一に溶ける溶解との直接性が最も低い2を選ぶ。ただしコロイド溶液という慣用語や高分子の溶解もあり、溶解とコロイドが一切無関係という意味ではない。

\Needspace{4\baselineskip}\paragraph{選択肢の検討（査読待ち初稿）}\mbox{}\par

\noindent\textbf{1：適切} エマルション。互いに混ざりにくい液体の一方が微小な液滴として他方へ分散した乳濁液である。\par

\noindent\textbf{2：条件付きで適切} 溶解。通常の小分子・イオンの溶解はコロイドの粒子分散とは区別する。高分子や会合体を含む系まで語だけで排除できない。\par

\noindent\textbf{3：適切} チンダル現象。分散粒子が光を散乱し、光の経路が見える現象である。\par

\noindent\textbf{4：適切} マヨネーズ。油滴が水性の連続相に分散し、卵黄成分などが乳化を安定化する例である。\par

\noindent\textbf{5：条件付きで適切} せっけん水。石けんは親水部と疎水部を持ち、適切な温度・濃度でミセル等の会合体を作る。全濃度の水溶液が同じコロイド状態という意味ではない。\par

\paragraph{再利用する考え方} 物質名より、何がどの尺度で分散しているかを読む。乳化した油滴と、石けん分子だけが集まったミセルを区別する。

\paragraph{共通理論} \hyperref[theory:phase_sep:colloid]{第\ref*{theory:phase_sep:colloid}節}, \hyperref[theory:phase_sep:micelle]{第\ref*{theory:phase_sep:micelle}節}

\par 参考：\cite{openstax_colloids}。

\par 難易度 2（校準初稿）、優先度未評価。

\section{真空排気・低温凝縮吸着}

\subsection{平成30年度 問8}
\label{q:h30-reg-applied-08}

\noindent\textbf{原典：} sources/H30/17\_applied\_science.pdf、PDF物理3頁。

\textbf{公式正答：} 1（single）。

\paragraph{問題の要旨} 気体の低温凝縮又は吸着を利用する真空ポンプを識別する。

\Needspace{4\baselineskip}\paragraph{確認した条件・記述（独自要約）}\mbox{}\par

\noindent\textbf{確認範囲：} 問題幹、条件、否定、数値、単位、全5肢を原典画像で照合した。\par

\noindent\textbf{成立条件：} 気体の凝縮又は吸着を利用する装置。拡散ポンプの油蒸気の循環と、排気対象を冷面へ捕捉する原理を区別する。\par

\paragraph{問われる核心} 気体の低温凝縮又は吸着を利用する真空ポンプを識別する。

\paragraph{必要知識} 気体の凝縮又は吸着を利用する装置。拡散ポンプの油蒸気の循環と、排気対象を冷面へ捕捉する原理を区別する。

\paragraph{考え方・正答を選ぶ理由（査読待ち初稿）} クライオポンプは低温面で気体を凝縮させ、凝縮しにくい気体を低温の吸着材へ捕捉する。気体分子を系内で捕捉する原理を問うため1を選ぶ。他の4候補は運動量輸送又は機械的な排気を主体とする。

\Needspace{4\baselineskip}\paragraph{選択肢の検討（査読待ち初稿）}\mbox{}\par

\noindent\textbf{1：適切} クライオポンプ。低温面での凝縮と、活性炭等の吸着材への低温吸着で排気対象を捕捉する。\par

\noindent\textbf{2：不適切} ターボ分子ポンプ。高速回転翼が気体分子へ方向を持つ運動量を与えて排出する。\par

\noindent\textbf{3：不適切} 拡散ポンプ。噴出する作動油等の蒸気が排気対象の気体を運ぶ。作動流体の凝縮があっても、対象気体を凝縮捕捉する主原理ではない。\par

\noindent\textbf{4：不適切} ダイヤフラムポンプ。膜の往復変形と弁で容積を変え、気体を吸入・排出する。\par

\noindent\textbf{5：不適切} ルーツ型ポンプ。互いにかみ合うロータで気体を吸入口から排気側へ移送する。\par

\paragraph{再利用する考え方} 排気対象の捕捉と、作動流体の相変化を分ける。クライオポンプには捕捉容量があり、再生が必要になる。

\paragraph{共通理論} \hyperref[theory:phase_sep:vacuum]{第\ref*{theory:phase_sep:vacuum}節}

\par 参考：\cite{leybold_cryopump}。

\par 難易度 2（校準初稿）、優先度未評価。

\section{蒸留・気液平衡比・共沸}

\subsection{平成27年度 問15}
\label{q:h27-reg-applied-15}

\noindent\textbf{原典：} sources/H27/17\_applied\_science.pdf、PDF物理5頁。

\textbf{公式正答：} 1（single）。

\par\textbf{検討未確定：} 平衡定数の意味を原典は限定しない。気液平衡比にもK値があり、他肢より関係が薄いことの一意な根拠はAstra査読へ残す。

\paragraph{問題の要旨} 蒸留の装置・操作・気液平衡を表す5用語の直接性を比較する。

\Needspace{4\baselineskip}\paragraph{確認した条件・記述（独自要約）}\mbox{}\par

\noindent\textbf{確認範囲：} 問題幹、条件、否定、数値、単位、全5肢を原典画像で照合した。\par

\noindent\textbf{成立条件：} 最も関係が薄い用語を選ぶ。反応平衡定数と相分配の気液平衡比を区別する。\par

\paragraph{問われる核心} 蒸留の装置・操作・気液平衡を表す5用語の直接性を比較する。

\paragraph{必要知識} 最も関係が薄い用語を選ぶ。反応平衡定数と相分配の気液平衡比を区別する。

\paragraph{考え方・正答を選ぶ理由（査読待ち初稿）} 理論段数、相対揮発度、還流、充填塔は精留の設計又は操作に直接現れる。公式は平衡定数の1を選ぶ。化学反応の平衡定数として読めばこの選択になるが、気液平衡のK値も蒸留に使用する。平衡という概念自体が蒸留と無関係という説明は成立しない。

\Needspace{4\baselineskip}\paragraph{選択肢の検討（査読待ち初稿）}\mbox{}\par

\noindent\textbf{1：解釈未解決} 平衡定数。化学反応の平衡定数という意味なら他候補より直接性が低いが、蒸留の気液平衡比 \(K_i=y_i/x_i\) もK値と呼ばれるため、関係がないとは断定できない。\par

\noindent\textbf{2：適切} 理論段数。各段で気液が平衡に達する理想化に必要な接触段数であり、精留塔の分離能力を表す。\par

\noindent\textbf{3：適切} 相対揮発度。2成分の気液平衡比の比であり、蒸留の分離の容易さを表す。\par

\noindent\textbf{4：適切} 還流。凝縮した塔頂液の一部を塔へ戻し、上昇蒸気と接触させる。\par

\noindent\textbf{5：適切} 充填塔。充填物の表面で液体と蒸気を接触させる装置であり、蒸留に使用する。\par

\paragraph{再利用する考え方} 蒸留は平衡を利用する分離操作である。\(K_i=y_i/x_i\) と反応の平衡定数の区別を、用語だけでなく定義から行う。

\paragraph{共通理論} \hyperref[theory:phase_sep:distillation]{第\ref*{theory:phase_sep:distillation}節}, \hyperref[theory:phase_sep:stages]{第\ref*{theory:phase_sep:stages}節}

\par 参考：\cite{nptel_distillation1}。

\paragraph{関連問題} \hyperref[q:r07-reg-applied-18]{R07-REG-APPLIED-18}

\par 難易度 2（校準初稿）、優先度未評価。

\subsection{令和元年度再試験 問10}
\label{q:r01-ret-applied-10}

\noindent\textbf{原典：} sources/R01\_retest/17\_applied\_science.pdf、PDF物理4頁。

\textbf{公式正答：} 4（single）。

\par\textbf{検討未確定：} 96\%の濃度基準を原典は指定しない。常圧の代表的な共沸値は質量基準約96\%と整合するが、厳密数値としての同定は追加査読事項に残す。

\paragraph{問題の要旨} 精留・水蒸気蒸留・減圧蒸留の目的と、共沸による通常精留の限界を比較する。

\Needspace{4\baselineskip}\paragraph{確認した条件・記述（独自要約）}\mbox{}\par

\noindent\textbf{確認範囲：} 問題幹、条件、否定、数値、単位、全5肢を原典画像で照合した。\par

\noindent\textbf{成立条件：} 外圧 \(1.013\times10^5\,\mathrm{Pa}\)。原典のエタノール96\%には質量・体積・モルの基準が明示されない。共沸の数値の解釈と、通常精留不能の原理を分ける。\par

\paragraph{問われる核心} 精留・水蒸気蒸留・減圧蒸留の目的と、共沸による通常精留の限界を比較する。

\paragraph{必要知識} 外圧 \(1.013\times10^5\,\mathrm{Pa}\)。原典のエタノール96\%には質量・体積・モルの基準が明示されない。共沸の数値の解釈と、通常精留不能の原理を分ける。

\paragraph{考え方・正答を選ぶ理由（査読待ち初稿）} 共沸組成では気相組成と液相組成が一致し、段数や還流を増してもその圧力での通常精留による分離は進まない。完全に水と分離できるとする4が不適切である。減圧は沸点を下げる操作、水蒸気蒸留はほぼ不混和な物質を水蒸気と留出させる操作であり、役割が異なる。

\Needspace{4\baselineskip}\paragraph{選択肢の検討（査読待ち初稿）}\mbox{}\par

\noindent\textbf{1：適切} より揮発性が高い成分は蒸気、低い成分は液に濃縮。同じ状態で気相への移りやすさの異なる成分を分ける。共沸点ではその差が失われるため一般操作の原理として読む。\par

\noindent\textbf{2：適切} 充填塔や棚段塔で還流・向流接触・蒸発と分縮を反復。上昇する蒸気と下降する液を接触させ、多段の分離に相当する効果を得る。\par

\noindent\textbf{3：適切} ほぼ混和しない物質へ水蒸気を吹き込み揮発成分を留出。水と対象の分圧が全圧へ寄与し、両者が共存する条件で水蒸気と共に成分を運ぶ。\par

\noindent\textbf{4：不適切} 約96\%のエタノール水溶液の共沸も通常の精留で完全分離できる。指定圧力の共沸では気液の組成が同じになり、通常の精留による濃縮が進まない。\par

\noindent\textbf{5：適切} 減圧蒸留は沸点を下げ熱分解を避けるのに適する。低い外圧で沸騰でき、熱に弱い成分の分離に適する。ただし減圧だけで全共沸が必ず消えるわけではない。\par

\paragraph{再利用する考え方} 分離装置の能力を増すことと、平衡曲線そのものを変えることは異なる。濃度パーセントは基準を確認し、96\%をモル分率と決め付けない。

\paragraph{共通理論} \hyperref[theory:phase_sep:distillation]{第\ref*{theory:phase_sep:distillation}節}, \hyperref[theory:phase_sep:azeotrope]{第\ref*{theory:phase_sep:azeotrope}節}, \hyperref[theory:phase_sep:steam]{第\ref*{theory:phase_sep:steam}節}

\par 参考：\cite{nptel_distillation1,nptel_azeotrope}。

\paragraph{関連問題} \hyperref[q:r06-reg-applied-15]{R06-REG-APPLIED-15}

\par 難易度 2（校準初稿）、優先度未評価。

\subsection{令和6年度 問15}
\label{q:r06-reg-applied-15}

\noindent\textbf{原典：} sources/R06/17\_applied\_science.pdf、PDF物理5頁。

\textbf{公式正答：} 1（single）。

\par\textbf{検討未確定：} 96\%の濃度基準を原典は指定しない。常圧の代表的な共沸値は質量基準約96\%と整合するが、厳密数値としての同定は追加査読事項に残す。

\paragraph{問題の要旨} 精留・水蒸気蒸留・減圧蒸留の目的と、共沸による通常精留の限界を比較する。

\Needspace{4\baselineskip}\paragraph{確認した条件・記述（独自要約）}\mbox{}\par

\noindent\textbf{確認範囲：} 問題幹、条件、否定、数値、単位、全5肢を原典画像で照合した。\par

\noindent\textbf{成立条件：} 外圧1気圧。原典のエタノール96\%には質量・体積・モルの基準が明示されない。共沸の数値の解釈と、通常精留不能の原理を分ける。\par

\paragraph{問われる核心} 精留・水蒸気蒸留・減圧蒸留の目的と、共沸による通常精留の限界を比較する。

\paragraph{必要知識} 外圧1気圧。原典のエタノール96\%には質量・体積・モルの基準が明示されない。共沸の数値の解釈と、通常精留不能の原理を分ける。

\paragraph{考え方・正答を選ぶ理由（査読待ち初稿）} 共沸組成では気相組成と液相組成が一致し、段数や還流を増してもその圧力での通常精留による分離は進まない。完全に水と分離できるとする1が不適切である。減圧は沸点を下げる操作、水蒸気蒸留はほぼ不混和な物質を水蒸気と留出させる操作であり、役割が異なる。

\Needspace{4\baselineskip}\paragraph{選択肢の検討（査読待ち初稿）}\mbox{}\par

\noindent\textbf{1：不適切} 約96\%のエタノール水溶液の共沸も通常の精留で完全分離できる。指定圧力の共沸では気液の組成が同じになり、通常の精留による濃縮が進まない。\par

\noindent\textbf{2：適切} 減圧蒸留は沸点を下げ熱分解を避けるのに適する。低い外圧で沸騰でき、熱に弱い成分の分離に適する。ただし減圧だけで全共沸が必ず消えるわけではない。\par

\noindent\textbf{3：適切} より揮発性が高い成分は蒸気、低い成分は液に濃縮。同じ状態で気相への移りやすさの異なる成分を分ける。共沸点ではその差が失われるため一般操作の原理として読む。\par

\noindent\textbf{4：適切} ほぼ混和しない物質へ水蒸気を吹き込み揮発成分を留出。水と対象の分圧が全圧へ寄与し、両者が共存する条件で水蒸気と共に成分を運ぶ。\par

\noindent\textbf{5：適切} 充填塔や棚段塔で還流・向流接触・蒸発と分縮を反復。上昇する蒸気と下降する液を接触させ、多段の分離に相当する効果を得る。\par

\paragraph{再利用する考え方} 分離装置の能力を増すことと、平衡曲線そのものを変えることは異なる。濃度パーセントは基準を確認し、96\%をモル分率と決め付けない。

\paragraph{共通理論} \hyperref[theory:phase_sep:distillation]{第\ref*{theory:phase_sep:distillation}節}, \hyperref[theory:phase_sep:azeotrope]{第\ref*{theory:phase_sep:azeotrope}節}, \hyperref[theory:phase_sep:steam]{第\ref*{theory:phase_sep:steam}節}

\par 参考：\cite{nptel_distillation1,nptel_azeotrope}。

\paragraph{関連問題} \hyperref[q:r01-ret-applied-10]{R01-RET-APPLIED-10}

\par 難易度 2（校準初稿）、優先度未評価。

\subsection{令和7年度 問18}
\label{q:r07-reg-applied-18}

\noindent\textbf{原典：} sources/R07/17\_applied\_science.pdf、PDF物理6頁。

\textbf{公式正答：} 4（single）。

\par\textbf{検討未確定：} 平衡定数の意味を原典は限定しない。気液平衡比にもK値があり、他肢より関係が薄いことの一意な根拠はAstra査読へ残す。

\paragraph{問題の要旨} 蒸留の装置・操作・気液平衡を表す5用語の直接性を比較する。

\Needspace{4\baselineskip}\paragraph{確認した条件・記述（独自要約）}\mbox{}\par

\noindent\textbf{確認範囲：} 問題幹、条件、否定、数値、単位、全5肢を原典画像で照合した。\par

\noindent\textbf{成立条件：} 最も関係が薄い用語を選ぶ。反応平衡定数と相分配の気液平衡比を区別する。\par

\paragraph{問われる核心} 蒸留の装置・操作・気液平衡を表す5用語の直接性を比較する。

\paragraph{必要知識} 最も関係が薄い用語を選ぶ。反応平衡定数と相分配の気液平衡比を区別する。

\paragraph{考え方・正答を選ぶ理由（査読待ち初稿）} 理論段数、相対揮発度、還流、充填塔は精留の設計又は操作に直接現れる。公式は平衡定数の4を選ぶ。化学反応の平衡定数として読めばこの選択になるが、気液平衡のK値も蒸留に使用する。平衡という概念自体が蒸留と無関係という説明は成立しない。

\Needspace{4\baselineskip}\paragraph{選択肢の検討（査読待ち初稿）}\mbox{}\par

\noindent\textbf{1：適切} 還流。凝縮した塔頂液の一部を塔へ戻し、上昇蒸気と接触させる。\par

\noindent\textbf{2：適切} 充填塔。充填物の表面で液体と蒸気を接触させる装置であり、蒸留に使用する。\par

\noindent\textbf{3：適切} 相対揮発度。2成分の気液平衡比の比であり、蒸留の分離の容易さを表す。\par

\noindent\textbf{4：解釈未解決} 平衡定数。化学反応の平衡定数という意味なら他候補より直接性が低いが、蒸留の気液平衡比 \(K_i=y_i/x_i\) もK値と呼ばれるため、関係がないとは断定できない。\par

\noindent\textbf{5：適切} 理論段数。各段で気液が平衡に達する理想化に必要な接触段数であり、精留塔の分離能力を表す。\par

\paragraph{再利用する考え方} 蒸留は平衡を利用する分離操作である。\(K_i=y_i/x_i\) と反応の平衡定数の区別を、用語だけでなく定義から行う。

\paragraph{共通理論} \hyperref[theory:phase_sep:distillation]{第\ref*{theory:phase_sep:distillation}節}, \hyperref[theory:phase_sep:stages]{第\ref*{theory:phase_sep:stages}節}

\par 参考：\cite{nptel_distillation1}。

\paragraph{関連問題} \hyperref[q:h27-reg-applied-15]{H27-REG-APPLIED-15}

\par 難易度 2（校準初稿）、優先度未評価。

\section{電池電位・反応自由エネルギー}

\subsection{令和4年度 問15}
\label{q:r04-reg-applied-15}

\noindent\textbf{原典：} sources/R04/17\_applied\_science.pdf、PDF物理5頁。

\textbf{公式正答：} 3（single）。

\paragraph{問題の要旨} 電池の可逆起電力と反応ギブズエネルギー、平衡、標準電極電位の関係を判定する。

\Needspace{4\baselineskip}\paragraph{確認した条件・記述（独自要約）}\mbox{}\par

\noindent\textbf{確認範囲：} 問題幹、数値、式、否定、全5候補を原典画像で照合した。\par

\noindent\textbf{成立条件：} 可逆・開回路の起電力を扱う。標準電極電位は共通の還元電位表で比較し、全反応の平衡と個別界面の平衡を分ける。\par

\paragraph{問われる核心} 電池の可逆起電力と反応ギブズエネルギー、平衡、標準電極電位の関係を判定する。

\paragraph{必要知識} 可逆・開回路の起電力を扱う。標準電極電位は共通の還元電位表で比較し、全反応の平衡と個別界面の平衡を分ける。

\paragraph{考え方・正答を選ぶ理由（査読待ち初稿）} 反応1 molの進行で外部へ運ばれる電荷はnFであり、可逆に取り出す電気仕事の最大値は \(-\Delta_rG\) である。したがって \(\Delta_rG=-nFE\) であり、無関係とする3が不適切である。全電池反応の化学平衡ではE=\allowbreak{}0となるが、個々の電極が局所平衡の開回路電池で起電力が0とは限らない。

\Needspace{4\baselineskip}\paragraph{選択肢の検討（査読待ち初稿）}\mbox{}\par

\noindent\textbf{1：適切} 酸化電極をアノード、還元電極をカソードと呼ぶ。反応による定義である。放電電池と電気分解では正負の対応が異なっても、酸化・還元による名称は変わらない。\par

\noindent\textbf{2：適切} 電池反応全体が化学平衡なら電位差0。全反応の \(\Delta_rG=0\) から \(E=0\)。個別電極界面が局所平衡である開回路状態だけとは区別する。\par

\noindent\textbf{3：不適切} 電池起電力は反応ギブズエネルギーと無関係。可逆起電力は \(\Delta_rG=-nFE\) で結ばれる。nは釣合い全反応で移動する電子のmol数である。\par

\noindent\textbf{4：適切} 標準水素電極を0とする基準。標準電極電位はこの参照電極に対する相対値として表す。単独電極の絶対電位を直接測定する意味ではない。\par

\noindent\textbf{5：適切} 標準電池電位は両標準電極電位の差。双方を還元電位として \(E^\circ=E^\circ_{\mathrm{cathode}}-E^\circ_{\mathrm{anode}}\)。反応の係数倍で電位を倍にしない。\par

\paragraph{再利用する考え方} 起電力は示強量である。反応式を倍にすると反応エネルギーと電子数が共に倍になり、Eは変わらない。

\paragraph{共通理論} \hyperref[theory:electrochem:gibbs]{第\ref*{theory:electrochem:gibbs}節}, \hyperref[theory:electrochem:nernst]{第\ref*{theory:electrochem:nernst}節}, \hyperref[theory:electrochem:electrodes]{第\ref*{theory:electrochem:electrodes}節}

\par 参考：\cite{openstax_cell_gibbs,mit_redox2014}。

\paragraph{関連問題} \hyperref[q:r07-reg-applied-10]{R07-REG-APPLIED-10}

\par 難易度 2（校準初稿）、優先度未評価。

\subsection{令和7年度 問10}
\label{q:r07-reg-applied-10}

\noindent\textbf{原典：} sources/R07/17\_applied\_science.pdf、PDF物理4頁。

\textbf{公式正答：} 1（single）。

\paragraph{問題の要旨} 電池の可逆起電力と反応ギブズエネルギー、平衡、標準電極電位の関係を判定する。

\Needspace{4\baselineskip}\paragraph{確認した条件・記述（独自要約）}\mbox{}\par

\noindent\textbf{確認範囲：} 問題幹、数値、式、否定、全5候補を原典画像で照合した。\par

\noindent\textbf{成立条件：} 可逆・開回路の起電力を扱う。標準電極電位は共通の還元電位表で比較し、全反応の平衡と個別界面の平衡を分ける。\par

\paragraph{問われる核心} 電池の可逆起電力と反応ギブズエネルギー、平衡、標準電極電位の関係を判定する。

\paragraph{必要知識} 可逆・開回路の起電力を扱う。標準電極電位は共通の還元電位表で比較し、全反応の平衡と個別界面の平衡を分ける。

\paragraph{考え方・正答を選ぶ理由（査読待ち初稿）} 反応1 molの進行で外部へ運ばれる電荷はnFであり、可逆に取り出す電気仕事の最大値は \(-\Delta_rG\) である。したがって \(\Delta_rG=-nFE\) であり、無関係とする1が不適切である。全電池反応の化学平衡ではE=\allowbreak{}0となるが、個々の電極が局所平衡の開回路電池で起電力が0とは限らない。

\Needspace{4\baselineskip}\paragraph{選択肢の検討（査読待ち初稿）}\mbox{}\par

\noindent\textbf{1：不適切} 電池起電力は反応ギブズエネルギーと無関係。可逆起電力は \(\Delta_rG=-nFE\) で結ばれる。nは釣合い全反応で移動する電子のmol数である。\par

\noindent\textbf{2：適切} 電池反応全体が化学平衡なら電位差0。全反応の \(\Delta_rG=0\) から \(E=0\)。個別電極界面が局所平衡である開回路状態だけとは区別する。\par

\noindent\textbf{3：適切} 酸化電極をアノード、還元電極をカソードと呼ぶ。反応による定義である。放電電池と電気分解では正負の対応が異なっても、酸化・還元による名称は変わらない。\par

\noindent\textbf{4：適切} 標準水素電極を0とする基準。標準電極電位はこの参照電極に対する相対値として表す。単独電極の絶対電位を直接測定する意味ではない。\par

\noindent\textbf{5：適切} 標準電池電位は両標準電極電位の差。双方を還元電位として \(E^\circ=E^\circ_{\mathrm{cathode}}-E^\circ_{\mathrm{anode}}\)。反応の係数倍で電位を倍にしない。\par

\paragraph{再利用する考え方} 起電力は示強量である。反応式を倍にすると反応エネルギーと電子数が共に倍になり、Eは変わらない。

\paragraph{共通理論} \hyperref[theory:electrochem:gibbs]{第\ref*{theory:electrochem:gibbs}節}, \hyperref[theory:electrochem:nernst]{第\ref*{theory:electrochem:nernst}節}, \hyperref[theory:electrochem:electrodes]{第\ref*{theory:electrochem:electrodes}節}

\par 参考：\cite{openstax_cell_gibbs,mit_redox2014}。

\paragraph{関連問題} \hyperref[q:r04-reg-applied-15]{R04-REG-APPLIED-15}

\par 難易度 2（校準初稿）、優先度未評価。

\section{ファラデーの法則・電気分解}

\subsection{平成25年度 問12}
\label{q:h25-reg-applied-12}

\noindent\textbf{原典：} sources/H25/17\_applied\_science.pdf、PDF物理5頁。

\textbf{公式正答：} 4（single）。

\paragraph{問題の要旨} 二価の銅イオンの析出量から必要電荷と通電時間を求める。

\Needspace{4\baselineskip}\paragraph{確認した条件・記述（独自要約）}\mbox{}\par

\noindent\textbf{確認範囲：} 問題幹、数値、式、否定、全5候補を原典画像で照合した。\par

\noindent\textbf{成立条件：} 電流5.36 A、Cu 1.27 g、モル質量63.5 g/mol、F=\allowbreak{}96500 C/mol。全電流が電気分解に使われるという原典指定により電流効率1とする。\par

\paragraph{問われる核心} 二価の銅イオンの析出量から必要電荷と通電時間を求める。

\paragraph{必要知識} 電流5.36 A、Cu 1.27 g、モル質量63.5 g/mol、F=\allowbreak{}96500 C/mol。全電流が電気分解に使われるという原典指定により電流効率1とする。

\paragraph{考え方・正答を選ぶ理由（査読待ち初稿）} 析出した銅は \(1.27/63.5=0.0200\,\mathrm{mol}\)。\(\mathrm{Cu^{2+}}+2e^-\to\mathrm{Cu}\) に必要な電子は0.0400 molである。\(Q=n_eF=0.0400\times96500=3860\,\mathrm C\)、\(t=Q/I=3860/5.36=720.15\,\mathrm s\) である。60で割ると約12.0分となり4を選ぶ。

\Needspace{4\baselineskip}\paragraph{選択肢の検討（査読待ち初稿）}\mbox{}\par

\noindent\textbf{1：不適切} 2分。必要電荷3860 Cに対し5.36×120=\allowbreak{}643.2 Cしか流れない。\par

\noindent\textbf{2：不適切} 6分。1929.6 Cであり、Cu2+に2電子が必要なことを無視するとこの時間に近付く。\par

\noindent\textbf{3：不適切} 10分。3216 Cであり、指定析出量に必要な3860 Cより不足する。\par

\noindent\textbf{4：適切} 12分。3859.2 Cとなり丸めた必要電荷と一致する。\par

\noindent\textbf{5：不適切} 20分。6432 Cとなり、指定析出量に必要な電荷より大きい。\par

\paragraph{再利用する考え方} 物質量→電子量→電荷→秒→分の順に単位をつないで計算する。Fを銅1 molの電荷とせず電子1 molの電荷として使う。

\paragraph{共通理論} \hyperref[theory:electrochem:faraday]{第\ref*{theory:electrochem:faraday}節}, \hyperref[theory:electrochem:copper]{第\ref*{theory:electrochem:copper}節}

\par 参考：\cite{mit_redox2014}。

\paragraph{関連問題} \hyperref[q:h26-reg-applied-13]{H26-REG-APPLIED-13}

\par 難易度 2（校準初稿）、優先度未評価。

\section{反応熱・ヘスの法則}

\subsection{平成24年度 問14}
\label{q:h24-reg-applied-14}

\noindent\textbf{原典：} sources/H24/17\_applied\_science.pdf、PDF物理5頁。

\textbf{公式正答：} 4（single）。

\paragraph{問題の要旨} 熱化学方程式を組み合わせ、指定量の液体ベンゼンの完全燃焼熱を求める。

\Needspace{4\baselineskip}\paragraph{確認した条件・記述（独自要約）}\mbox{}\par

\noindent\textbf{確認範囲：} 問題幹、数値、式、否定、全5候補を原典画像で照合した。\par

\noindent\textbf{成立条件：} 原典の水は液体、炭素は黒鉛、ベンゼンは液体。式に付く熱量の符号を反応エンタルピーへ変換し、指定物質量を最後に掛ける。\par

\paragraph{問われる核心} 熱化学方程式を組み合わせ、指定量の液体ベンゼンの完全燃焼熱を求める。

\paragraph{必要知識} 原典の水は液体、炭素は黒鉛、ベンゼンは液体。式に付く熱量の符号を反応エンタルピーへ変換し、指定物質量を最後に掛ける。

\paragraph{考え方・正答を選ぶ理由（査読待ち初稿）} ベンゼンの生成式の右辺にある−49 kJは吸熱を表し、生成エンタルピーは+49 kJ/molである。燃焼式は \(\mathrm{C_6H_6(l)}+(15/2)\mathrm{O_2(g)}\to6\mathrm{CO_2(g)}+3\mathrm{H_2O(l)}\)。ヘスの法則で \(\Delta H=6(-394)+3(-286)-49=-3271\,\mathrm{kJ\,mol^{-1}}\) となる。1 molの放出熱量は3271 kJであり、最も近い3300 kJの4を選ぶ。

\Needspace{4\baselineskip}\paragraph{選択肢の検討（査読待ち初稿）}\mbox{}\par

\noindent\textbf{1：不適切} 900 kJ。生成物のCO2六つと水三つの寄与を合計した値より大きく不足する。\par

\noindent\textbf{2：不適切} 1700 kJ。炭素と水素の両方の酸化による発熱を正しく合計した値に達しない。\par

\noindent\textbf{3：不適切} 2500 kJ。生成物と反応物のエンタルピー差から求めた値より小さい。\par

\noindent\textbf{4：適切} 3300 kJ。原典指定の量の放出熱量に最も近い。\par

\noindent\textbf{5：不適切} 4100 kJ。原典の反応係数と指定物質量から求めた値より大きい。\par

\paragraph{再利用する考え方} 反応を逆にするとエンタルピーの符号が逆転し、係数倍すると熱量も同じ倍率になる。放出熱量の正値と負の反応エンタルピーを区別する。

\paragraph{共通理論} \hyperref[theory:electrochem:hess]{第\ref*{theory:electrochem:hess}節}, \hyperref[theory:electrochem:heat_sign]{第\ref*{theory:electrochem:heat_sign}節}, \hyperref[theory:electrochem:combustion]{第\ref*{theory:electrochem:combustion}節}

\par 参考：\cite{openstax_hess}。

\paragraph{関連問題} \hyperref[q:h26-reg-applied-15]{H26-REG-APPLIED-15}, \hyperref[q:h28-reg-applied-12]{H28-REG-APPLIED-12}, \hyperref[q:r01-ret-applied-12]{R01-RET-APPLIED-12}

\par 難易度 2（校準初稿）、優先度未評価。

\subsection{平成28年度 問12}
\label{q:h28-reg-applied-12}

\noindent\textbf{原典：} sources/H28/17\_applied\_science.pdf、PDF物理5頁。

\textbf{公式正答：} 4（single）。

\paragraph{問題の要旨} 熱化学方程式を組み合わせ、指定量の液体ベンゼンの完全燃焼熱を求める。

\Needspace{4\baselineskip}\paragraph{確認した条件・記述（独自要約）}\mbox{}\par

\noindent\textbf{確認範囲：} 問題幹、数値、式、否定、全5候補を原典画像で照合した。\par

\noindent\textbf{成立条件：} 原典の水は液体、炭素は黒鉛、ベンゼンは液体。式に付く熱量の符号を反応エンタルピーへ変換し、指定物質量を最後に掛ける。\par

\paragraph{問われる核心} 熱化学方程式を組み合わせ、指定量の液体ベンゼンの完全燃焼熱を求める。

\paragraph{必要知識} 原典の水は液体、炭素は黒鉛、ベンゼンは液体。式に付く熱量の符号を反応エンタルピーへ変換し、指定物質量を最後に掛ける。

\paragraph{考え方・正答を選ぶ理由（査読待ち初稿）} ベンゼンの生成式の右辺にある−49 kJは吸熱を表し、生成エンタルピーは+49 kJ/molである。燃焼式は \(\mathrm{C_6H_6(l)}+(15/2)\mathrm{O_2(g)}\to6\mathrm{CO_2(g)}+3\mathrm{H_2O(l)}\)。ヘスの法則で \(\Delta H=6(-394)+3(-286)-49=-3271\,\mathrm{kJ\,mol^{-1}}\) となる。1 molの放出熱量は3271 kJであり、最も近い3300 kJの4を選ぶ。

\Needspace{4\baselineskip}\paragraph{選択肢の検討（査読待ち初稿）}\mbox{}\par

\noindent\textbf{1：不適切} 900 kJ。生成物のCO2六つと水三つの寄与を合計した値より大きく不足する。\par

\noindent\textbf{2：不適切} 1700 kJ。炭素と水素の両方の酸化による発熱を正しく合計した値に達しない。\par

\noindent\textbf{3：不適切} 2500 kJ。生成物と反応物のエンタルピー差から求めた値より小さい。\par

\noindent\textbf{4：適切} 3300 kJ。原典指定の量の放出熱量に最も近い。\par

\noindent\textbf{5：不適切} 4100 kJ。原典の反応係数と指定物質量から求めた値より大きい。\par

\paragraph{再利用する考え方} 反応を逆にするとエンタルピーの符号が逆転し、係数倍すると熱量も同じ倍率になる。放出熱量の正値と負の反応エンタルピーを区別する。

\paragraph{共通理論} \hyperref[theory:electrochem:hess]{第\ref*{theory:electrochem:hess}節}, \hyperref[theory:electrochem:heat_sign]{第\ref*{theory:electrochem:heat_sign}節}, \hyperref[theory:electrochem:combustion]{第\ref*{theory:electrochem:combustion}節}

\par 参考：\cite{openstax_hess}。

\paragraph{関連問題} \hyperref[q:h24-reg-applied-14]{H24-REG-APPLIED-14}, \hyperref[q:r01-ret-applied-12]{R01-RET-APPLIED-12}

\par 難易度 2（校準初稿）、優先度未評価。

\subsection{令和元年度再試験 問12}
\label{q:r01-ret-applied-12}

\noindent\textbf{原典：} sources/R01\_retest/17\_applied\_science.pdf、PDF物理5頁。

\textbf{公式正答：} 4（single）。

\paragraph{問題の要旨} 熱化学方程式を組み合わせ、指定量の液体ベンゼンの完全燃焼熱を求める。

\Needspace{4\baselineskip}\paragraph{確認した条件・記述（独自要約）}\mbox{}\par

\noindent\textbf{確認範囲：} 問題幹、数値、式、否定、全5候補を原典画像で照合した。\par

\noindent\textbf{成立条件：} 原典の水は液体、炭素は黒鉛、ベンゼンは液体。式に付く熱量の符号を反応エンタルピーへ変換し、指定物質量を最後に掛ける。\par

\paragraph{問われる核心} 熱化学方程式を組み合わせ、指定量の液体ベンゼンの完全燃焼熱を求める。

\paragraph{必要知識} 原典の水は液体、炭素は黒鉛、ベンゼンは液体。式に付く熱量の符号を反応エンタルピーへ変換し、指定物質量を最後に掛ける。

\paragraph{考え方・正答を選ぶ理由（査読待ち初稿）} ベンゼンの生成式の右辺にある−49 kJは吸熱を表し、生成エンタルピーは+49 kJ/molである。燃焼式は \(\mathrm{C_6H_6(l)}+(15/2)\mathrm{O_2(g)}\to6\mathrm{CO_2(g)}+3\mathrm{H_2O(l)}\)。ヘスの法則で \(\Delta H=6(-394)+3(-286)-49=-3271\,\mathrm{kJ\,mol^{-1}}\) となる。0.5 molでは放出熱量が \(0.5\times3271=1635.5\,\mathrm{kJ}\) なので1650 kJの4を選ぶ。

\Needspace{4\baselineskip}\paragraph{選択肢の検討（査読待ち初稿）}\mbox{}\par

\noindent\textbf{1：不適切} 450 kJ。生成物のCO2六つと水三つの寄与を合計した値より大きく不足する。\par

\noindent\textbf{2：不適切} 850 kJ。炭素と水素の両方の酸化による発熱を正しく合計した値に達しない。\par

\noindent\textbf{3：不適切} 1250 kJ。生成物と反応物のエンタルピー差から求めた値より小さい。\par

\noindent\textbf{4：適切} 1650 kJ。原典指定の量の放出熱量に最も近い。\par

\noindent\textbf{5：不適切} 2050 kJ。原典の反応係数と指定物質量から求めた値より大きい。\par

\paragraph{再利用する考え方} 反応を逆にするとエンタルピーの符号が逆転し、係数倍すると熱量も同じ倍率になる。放出熱量の正値と負の反応エンタルピーを区別する。

\paragraph{共通理論} \hyperref[theory:electrochem:hess]{第\ref*{theory:electrochem:hess}節}, \hyperref[theory:electrochem:heat_sign]{第\ref*{theory:electrochem:heat_sign}節}, \hyperref[theory:electrochem:combustion]{第\ref*{theory:electrochem:combustion}節}

\par 参考：\cite{openstax_hess}。

\paragraph{関連問題} \hyperref[q:h24-reg-applied-14]{H24-REG-APPLIED-14}, \hyperref[q:h28-reg-applied-12]{H28-REG-APPLIED-12}

\par 難易度 2（校準初稿）、優先度未評価。

\section{鉛蓄電池・電解液質量収支}

\subsection{平成26年度 問13}
\label{q:h26-reg-applied-13}

\noindent\textbf{原典：} sources/H26/17\_applied\_science.pdf、PDF物理5頁。

\textbf{公式正答：} 4（single）。

\par\textbf{検討未確定：} 電流効率と外部の水損失を原典は明示しない。「使用後」を放電と解釈した主反応のみのモデルで公式4と一致する。充電や副反応を含む一般条件へ広げない。

\paragraph{問題の要旨} 鉛蓄電池の放電反応から、硫酸量と電解液量の両方を更新して濃度を求める。

\Needspace{4\baselineskip}\paragraph{確認した条件・記述（独自要約）}\mbox{}\par

\noindent\textbf{確認範囲：} 問題幹、数値、式、否定、全5候補を原典画像で照合した。\par

\noindent\textbf{成立条件：} 30質量\%の電解液1000 g、100 A、965 s、F=\allowbreak{}96500 C/mol。使用後を放電と読み、全電荷が主反応へ使われ、外部との液交換や蒸発がないモデルを用いる。\par

\paragraph{問われる核心} 鉛蓄電池の放電反応から、硫酸量と電解液量の両方を更新して濃度を求める。

\paragraph{必要知識} 30質量\%の電解液1000 g、100 A、965 s、F=\allowbreak{}96500 C/mol。使用後を放電と読み、全電荷が主反応へ使われ、外部との液交換や蒸発がないモデルを用いる。

\paragraph{考え方・正答を選ぶ理由（査読待ち初稿）} 放電の全反応は \(\mathrm{Pb+PbO_2+2H_2SO_4\to2PbSO_4+2H_2O}\) で、2 molの電子に対し硫酸2 molを消費し水2 molを生成する。通過電子は \(100\times965/96500=1.00\,\mathrm{mol}\)。硫酸98 gが減り水18 gが増える。硫酸は \(300-98=202\,\mathrm g\)、電解液全体は \(1000-98+18=920\,\mathrm g\) となる。質量百分率 \(100\times202/920=21.96\%\) に最も近い22\%の4を選ぶ。

\Needspace{4\baselineskip}\paragraph{選択肢の検討（査読待ち初稿）}\mbox{}\par

\noindent\textbf{1：不適切} 11\%。電子1 molに対する硫酸消費量1 molと水生成量1 molを使った質量収支に合わない。\par

\noindent\textbf{2：不適切} 16\%。硫酸の減少と水の生成を合わせた最終質量920 gに対する202 gの比ではない。\par

\noindent\textbf{3：不適切} 20\%。202/1000を約20\%とするのは電解液全質量を初期値のまま固定する誤りである。\par

\noindent\textbf{4：適切} 22\%。硫酸202 gを最終電解液920 gで割った約21.96\%に最も近い。\par

\noindent\textbf{5：不適切} 25\%。反応による硫酸消費98 gを反映した最終比率より大きい。\par

\paragraph{再利用する考え方} 電極上の固体PbSO4へ移る質量は液相の分母から除く。硫酸が水になるという言い換えだけで分母を1000 gのまま保たない。

\paragraph{共通理論} \hyperref[theory:electrochem:lead]{第\ref*{theory:electrochem:lead}節}, \hyperref[theory:electrochem:solution_mass]{第\ref*{theory:electrochem:solution_mass}節}, \hyperref[theory:electrochem:faraday]{第\ref*{theory:electrochem:faraday}節}

\par 参考：\cite{openstax_batteries,mit_redox2014}。

\paragraph{関連問題} \hyperref[q:h25-reg-applied-12]{H25-REG-APPLIED-12}

\par 難易度 3（校準初稿）、優先度未評価。

\section{リチウム電池・イオン挿入}

\subsection{平成24年度 問17}
\label{q:h24-reg-applied-17}

\noindent\textbf{原典：} sources/H24/17\_applied\_science.pdf、PDF物理6頁。

\textbf{公式正答：} 4（single）。

\par\textbf{検討未確定：} 「最も卑」は電極反応・溶媒・標準状態を原典が明記しない。通常の水溶液標準還元電位の比較として扱い、全溶媒系への全称とはしない。

\paragraph{問題の要旨} リチウムイオン電池の材料と、金属リチウム負極の課題を5下線部で判定する。

\Needspace{4\baselineskip}\paragraph{確認した条件・記述（独自要約）}\mbox{}\par

\noindent\textbf{確認範囲：} 問題幹、数値、式、否定、全5候補を原典画像で照合した。\par

\noindent\textbf{成立条件：} 原典の下線4はOCRの「自金」でなく「白金」である。標準電極電位は通常の水溶液系還元電位の基準、負極は放電時の役割として読む。\par

\paragraph{問われる核心} リチウムイオン電池の材料と、金属リチウム負極の課題を5下線部で判定する。

\paragraph{必要知識} 原典の下線4はOCRの「自金」でなく「白金」である。標準電極電位は通常の水溶液系還元電位の基準、負極は放電時の役割として読む。

\paragraph{考え方・正答を選ぶ理由（査読待ち初稿）} 金属リチウムの析出によるデンドライトは内部短絡を起こし得る。実用化で負極にはリチウムイオンを取り込む炭素系材料が採用され、白金という下線4が不適切である。初期の石油コークスと後の黒鉛を歴史上同一材料としない。正極のコバルト酸リチウムはLiCoO2である。

\Needspace{4\baselineskip}\paragraph{選択肢の検討（査読待ち初稿）}\mbox{}\par

\noindent\textbf{1：適切} 卑な。通常の標準還元電位でLi+/Liは非常に負である。金属リチウムを酸化させる高電圧電池の負極としての性質につながる。\par

\noindent\textbf{2：適切} 小さい。リチウムは金属元素で最も小さいモル質量を持つ。1 molあたりの移動電子数と合わせて質量当たり容量を考える。\par

\noindent\textbf{3：適切} 負極。放電時にリチウムを酸化させ、電子を外部へ供給する側である。\par

\noindent\textbf{4：不適切} 白金。この実用化経路の負極は炭素系材料であり、白金ではない。白金の不活性電極としての別の用途と混同しない。\par

\noindent\textbf{5：適切} コバルト酸リチウム。LiCoO2はリチウムイオンの出入りを利用する実用正極材料である。\par

\paragraph{再利用する考え方} 高い標準電位差だけで実用電池は決まらない。可逆にイオンを取り込むホスト、電解液、界面、短絡の抑制を合わせて読む。

\paragraph{共通理論} \hyperref[theory:electrochem:lithium]{第\ref*{theory:electrochem:lithium}節}, \hyperref[theory:electrochem:electrodes]{第\ref*{theory:electrochem:electrodes}節}

\par 参考：\cite{nobel_lithium2019,openstax_batteries}。

\par 難易度 2（校準初稿）、優先度未評価。

\section{相状態・臨界・実在気体}

\subsection{平成26年度 問10}
\label{q:h26-reg-applied-10}

\noindent\textbf{原典：} sources/H26/17\_applied\_science.pdf、PDF物理4頁。

\textbf{公式正答：} 3（single）。

\par\textbf{検討未確定：} 熱容量の差を仕事だけで説明する原典肢には理想気体の指定がない。純水の相平衡温度は丸めた常用値である。

\paragraph{問題の要旨} 相平衡の化学ポテンシャル、臨界点、実在気体と定圧熱容量の条件を比較する。

\Needspace{4\baselineskip}\paragraph{確認した条件・記述（独自要約）}\mbox{}\par

\noindent\textbf{確認範囲：} 問題幹、条件、否定、数値、単位、全5肢を原典画像で照合した。\par

\noindent\textbf{成立条件：} 標準状態でのCO2と窒素の比較。純物質の相平衡、圧縮による液化、昇華、同じ昇温の熱容量比較。\par

\paragraph{問われる核心} 相平衡の化学ポテンシャル、臨界点、実在気体と定圧熱容量の条件を比較する。

\paragraph{必要知識} 標準状態でのCO2と窒素の比較。純物質の相平衡、圧縮による液化、昇華、同じ昇温の熱容量比較。

\paragraph{考え方・正答を選ぶ理由（査読待ち初稿）} 相平衡は相ごとの化学ポテンシャルの一致であり、臨界点では気液共存が終わる。CO2の方が窒素より理想性のずれが小さいという比較が逆であるため、3を選ぶ。熱容量差を外部仕事だけで説明する肢は理想気体を想定する必要がある。

\Needspace{4\baselineskip}\paragraph{選択肢の検討（査読待ち初稿）}\mbox{}\par

\noindent\textbf{1：条件付きで適切} 純水のモルGは1気圧100°Cで気液、0°Cで固液が等しい。純物質の相平衡条件である。原典の0°C・100°Cは常用の丸め値であり、精密な共存温度の定義ではない。\par

\noindent\textbf{2：適切} 臨界点では気液の区別が消え、臨界温度以上で圧縮だけでは通常の液化ができない。臨界点は気液共存曲線の終端である。超臨界の流体を高密度にすることと、気液界面を持つ液化は異なる。\par

\noindent\textbf{3：不適切} 標準状態で二酸化炭素は窒素より理想気体からのずれが小さい。低圧の標準的条件ではCO2の相互作用による非理想性が窒素より大きい。理想性の比較は同じ温度・圧力で行い、気体名だけで全状態へ一般化しない。\par

\noindent\textbf{4：条件付きで適切} 定圧加熱が定積より多く熱を要する差は外部仕事による。理想気体の同じ初期状態と同じ昇温なら \(U=U(T)\) であり差は膨張仕事となる。実在気体では最終体積の違いに伴う内部エネルギー差もあり、仕事だけという説明は一般化できない。\par

\noindent\textbf{5：適切} 固体が液相を経ず気体に昇華することがある。固相と気相が直接関わる相変化であり、氷も温度・圧力条件に応じて昇華する。\par

\paragraph{再利用する考え方} 相平衡の等式と、実在気体への近似を区別する。定圧・定積で同じ最終温度でも最終体積は異なり、実在気体の内部エネルギーは一般に同じにならない。

\paragraph{共通理論} \hyperref[theory:phase_sep:chemical_potential]{第\ref*{theory:phase_sep:chemical_potential}節}, \hyperref[theory:phase_sep:critical]{第\ref*{theory:phase_sep:critical}節}, \hyperref[theory:phase_sep:realgas]{第\ref*{theory:phase_sep:realgas}節}, \hyperref[theory:phase_sep:capacity]{第\ref*{theory:phase_sep:capacity}節}

\par 参考：\cite{iupac_greenbook2012,mit_thermo_firstlaw2008}。

\paragraph{関連問題} \hyperref[q:r01-reg-applied-16]{R01-REG-APPLIED-16}

\par 難易度 2（校準初稿）、優先度未評価。

\subsection{令和元年度 問16}
\label{q:r01-reg-applied-16}

\noindent\textbf{原典：} sources/R01/17\_applied\_science.pdf、PDF物理5頁。

\textbf{公式正答：} 4（single）。

\par\textbf{検討未確定：} 熱容量の差を仕事だけで説明する原典肢には理想気体の指定がない。

\paragraph{問題の要旨} 相平衡の化学ポテンシャル、臨界点、実在気体と定圧熱容量の条件を比較する。

\Needspace{4\baselineskip}\paragraph{確認した条件・記述（独自要約）}\mbox{}\par

\noindent\textbf{確認範囲：} 問題幹、条件、否定、数値、単位、全5肢を原典画像で照合した。\par

\noindent\textbf{成立条件：} 標準状態でのCO2と窒素の比較。純物質の相平衡、圧縮による液化、昇華、同じ昇温の熱容量比較。\par

\paragraph{問われる核心} 相平衡の化学ポテンシャル、臨界点、実在気体と定圧熱容量の条件を比較する。

\paragraph{必要知識} 標準状態でのCO2と窒素の比較。純物質の相平衡、圧縮による液化、昇華、同じ昇温の熱容量比較。

\paragraph{考え方・正答を選ぶ理由（査読待ち初稿）} 相平衡は相ごとの化学ポテンシャルの一致であり、臨界点では気液共存が終わる。CO2の方が窒素より理想性のずれが小さいという比較が逆であるため、4を選ぶ。熱容量差を外部仕事だけで説明する肢は理想気体を想定する必要がある。

\Needspace{4\baselineskip}\paragraph{選択肢の検討（査読待ち初稿）}\mbox{}\par

\noindent\textbf{1：適切} 相平衡にある固体と液体の化学ポテンシャルは等しい。相間の物質移動でギブズエネルギーが低下しない条件である。純物質の化学ポテンシャルはモルギブズエネルギーである。\par

\noindent\textbf{2：適切} 臨界点では気液の区別が消え、臨界温度以上で圧縮だけでは通常の液化ができない。臨界点は気液共存曲線の終端である。超臨界の流体を高密度にすることと、気液界面を持つ液化は異なる。\par

\noindent\textbf{3：適切} 固体が液相を経ず気体に昇華することがある。固相と気相が直接関わる相変化であり、氷も温度・圧力条件に応じて昇華する。\par

\noindent\textbf{4：不適切} 標準状態で二酸化炭素は窒素より理想気体からのずれが小さい。低圧の標準的条件ではCO2の相互作用による非理想性が窒素より大きい。理想性の比較は同じ温度・圧力で行い、気体名だけで全状態へ一般化しない。\par

\noindent\textbf{5：条件付きで適切} 定圧加熱が定積より多く熱を要する差は外部仕事による。理想気体の同じ初期状態と同じ昇温なら \(U=U(T)\) であり差は膨張仕事となる。実在気体では最終体積の違いに伴う内部エネルギー差もあり、仕事だけという説明は一般化できない。\par

\paragraph{再利用する考え方} 相平衡の等式と、実在気体への近似を区別する。定圧・定積で同じ最終温度でも最終体積は異なり、実在気体の内部エネルギーは一般に同じにならない。

\paragraph{共通理論} \hyperref[theory:phase_sep:chemical_potential]{第\ref*{theory:phase_sep:chemical_potential}節}, \hyperref[theory:phase_sep:critical]{第\ref*{theory:phase_sep:critical}節}, \hyperref[theory:phase_sep:realgas]{第\ref*{theory:phase_sep:realgas}節}, \hyperref[theory:phase_sep:capacity]{第\ref*{theory:phase_sep:capacity}節}

\par 参考：\cite{iupac_greenbook2012,mit_thermo_firstlaw2008}。

\paragraph{関連問題} \hyperref[q:h26-reg-applied-10]{H26-REG-APPLIED-10}

\par 難易度 2（校準初稿）、優先度未評価。

\subsection{令和3年度 問15}
\label{q:r03-reg-applied-15}

\noindent\textbf{原典：} sources/R03/17\_applied\_science.pdf、PDF物理6頁。

\textbf{公式正答：} 3（single）。

\par\textbf{検討未確定：} 気体の種類を原典は指定しない。冷却の全称には理想気体の反例があり、ギブズエネルギーの判定条件には温度一定等が省略される。

\paragraph{問題の要旨} 外圧0の断熱膨張での保存量を確認し、気体が必ず冷えるかを検討する。

\Needspace{4\baselineskip}\paragraph{確認した条件・記述（独自要約）}\mbox{}\par

\noindent\textbf{確認範囲：} 問題幹、条件、否定、数値、単位、全5肢を原典画像で照合した。\par

\noindent\textbf{成立条件：} 一定量の気体、外圧0、断熱、前後は平衡。外部仕事を行う断熱膨張と自由膨張を区別する。天秤の浮力は空気の密度と物体体積に依存する。\par

\paragraph{問われる核心} 外圧0の断熱膨張での保存量を確認し、気体が必ず冷えるかを検討する。

\paragraph{必要知識} 一定量の気体、外圧0、断熱、前後は平衡。外部仕事を行う断熱膨張と自由膨張を区別する。天秤の浮力は空気の密度と物体体積に依存する。

\paragraph{考え方・正答を選ぶ理由（査読待ち初稿）} 断熱で \(Q=0\)、外圧0で体積仕事 \(W=0\) なので \(\Delta U=0\) である。理想気体なら \(U=U(T)\) より温度は変わらない。実在気体では \((\partial T/\partial V)_U=[p-T(\partial p/\partial T)_V]/C_V\) の符号が状態方程式に依存する。必ず冷却するとする3は一般に成立しない。

\Needspace{4\baselineskip}\paragraph{選択肢の検討（査読待ち初稿）}\mbox{}\par

\noindent\textbf{1：適切} 三相が同時に平衡共存する点は三重点。一成分系の三相平衡を表す。相図上の一点であり臨界点とは異なる。\par

\noindent\textbf{2：条件付きで適切} ギブズエネルギーは定圧条件の自由エネルギー。圧力を自然変数とする \(G=U+pV-TS\) である。自発変化や極小の判定には温度も一定、閉じた系で非体積仕事がない等の条件が必要である。\par

\noindent\textbf{3：不適切} 真空への断熱膨張後は必ず冷却される。理想気体では温度不変が反例となる。実在気体でも内部エネルギー一定に沿う温度変化は状態方程式に依存する。\par

\noindent\textbf{4：適切} ファンデルワールス式は実在気体の体積と分子間相互作用を考慮。原典の実在気体の体積を、分子の有限サイズによる排除体積bの補正として説明する。引力はaの項で表す。\par

\noindent\textbf{5：適切} 天秤の測定では温湿度による空気密度の変化に留意。排除する空気質量による浮力が試料と分銅の見掛けの重量を変える。温湿度は空気密度を変える。\par

\paragraph{再利用する考え方} 断熱という条件だけでなく外圧と仕事を読む。自由膨張の内部エネルギー一定を、絞り流れのエンタルピー一定と混同しない。

\paragraph{共通理論} \hyperref[theory:phase_sep:free]{第\ref*{theory:phase_sep:free}節}, \hyperref[theory:phase_sep:gibbs]{第\ref*{theory:phase_sep:gibbs}節}, \hyperref[theory:phase_sep:realgas]{第\ref*{theory:phase_sep:realgas}節}, \hyperref[theory:phase_sep:buoyancy]{第\ref*{theory:phase_sep:buoyancy}節}

\par 参考：\cite{mit_thermo_firstlaw2008,iupac_greenbook2012}。

\par 難易度 2（校準初稿）、優先度未評価。

\subsection{令和5年度 問16}
\label{q:r05-reg-applied-16}

\noindent\textbf{原典：} sources/R05/17\_applied\_science.pdf、PDF物理6頁。

\textbf{公式正答：} 1（single）。

\par\textbf{検討未確定：} 平均密度・初期充填率が原典で指定されない。気相密度の単調上昇や任意容器の臨界到達を断定せず、臨界密度を持つ反例として公式1を説明する。

\paragraph{問題の要旨} 剛体密閉容器の平均密度を固定した加熱と、臨界点への到達条件を検討する。

\Needspace{4\baselineskip}\paragraph{確認した条件・記述（独自要約）}\mbox{}\par

\noindent\textbf{確認範囲：} 問題幹、条件、否定、数値、単位、全5肢を原典画像で照合した。\par

\noindent\textbf{成立条件：} 変形しない閉じた容器、加熱、平均密度一定。臨界密度322 kg/m3、臨界温度647.096 KはIAPWS-95の参照定数である。\par

\paragraph{問われる核心} 剛体密閉容器の平均密度を固定した加熱と、臨界点への到達条件を検討する。

\paragraph{必要知識} 変形しない閉じた容器、加熱、平均密度一定。臨界密度322 kg/m3、臨界温度647.096 KはIAPWS-95の参照定数である。

\paragraph{考え方・正答を選ぶ理由（査読待ち初稿）} 臨界点へ二相共存状態から近づくと液体と蒸気の密度差は0となる。剛体閉容器では平均密度 \(\bar\rho=M/V\) が一定であり、\(\bar\rho=\rho_c\) に選べば臨界点へ達する経路が可能である。したがって密度が一致することはないという1の全称は誤る。ただし任意の充填量で必ず臨界へ達するわけではなく、液相又は気相が先に消える場合もある。

\Needspace{4\baselineskip}\paragraph{選択肢の検討（査読待ち初稿）}\mbox{}\par

\noindent\textbf{1：不適切} 加熱した蒸気の密度は液体の密度に一致することがない。平均密度を臨界密度に設定した加熱が反例となる。任意密閉容器での一致を主張する必要はない。\par

\noindent\textbf{2：適切} ファンデルワールス式は分子間相互作用を考慮。分子間引力のa項を持つ。排除体積のb項も理想気体との差となる。\par

\noindent\textbf{3：条件付きで適切} ギブズエネルギーは定圧条件の自由エネルギー。圧力を自然変数とする \(G=U+pV-TS\) である。自発変化や極小の判定には温度も一定、閉じた系で非体積仕事がない等の条件が必要である。\par

\noindent\textbf{4：適切} 気液固が同時に平衡共存する点は三重点。一成分三相共存の定義である。\par

\noindent\textbf{5：適切} 気液平衡では化学ポテンシャルが等しい。相間の物質移動が全ギブズエネルギーを低下させない条件である。\par

\paragraph{再利用する考え方} 常に起こらないという全称は一つの可能な経路で反証できる。臨界点への加熱には温度だけでなく平均密度の条件が必要である。

\paragraph{共通理論} \hyperref[theory:phase_sep:critical]{第\ref*{theory:phase_sep:critical}節}, \hyperref[theory:phase_sep:isochore]{第\ref*{theory:phase_sep:isochore}節}, \hyperref[theory:phase_sep:gibbs]{第\ref*{theory:phase_sep:gibbs}節}

\par 参考：\cite{iapws95_2018}。

\par 難易度 2（校準初稿）、優先度未評価。

\section{反応商による化学平衡移動}

\subsection{平成28年度 問18}
\label{q:h28-reg-applied-18}

\noindent\textbf{原典：} sources/H28/17\_applied\_science.pdf、PDF物理7頁。

\textbf{公式正答：} 3（single）。

\paragraph{問題の要旨} 発熱するアンモニア合成の温度・圧力・窒素追加・触媒の効果を判定する。

\Needspace{4\baselineskip}\paragraph{確認した条件・記述（独自要約）}\mbox{}\par

\noindent\textbf{確認範囲：} 問題幹、否定、全5肢、数式、数値、単位を原典画像で確認した。\par

\noindent\textbf{成立条件：} \(\mathrm{N_2+3H_2\rightleftharpoons2NH_3}+92\,\mathrm{kJ}\)。窒素追加は温度・体積一定。気体の分圧を用いる理想気体近似。\par

\paragraph{問われる核心} 発熱するアンモニア合成の温度・圧力・窒素追加・触媒の効果を判定する。

\paragraph{必要知識} \(\mathrm{N_2+3H_2\rightleftharpoons2NH_3}+92\,\mathrm{kJ}\)。窒素追加は温度・体積一定。気体の分圧を用いる理想気体近似。

\paragraph{考え方・正答を選ぶ理由（査読待ち初稿）} 反応式は窒素1、水素3、アンモニア2であり、生成側に92 kJを持つ。窒素追加は定温定積で反応商を低下させる。低温は発熱方向、高圧は気体量の少ない生成方向を平衡上有利にする。触媒による平衡移動だけが誤りであり、3を選ぶ。

\Needspace{4\baselineskip}\paragraph{選択肢の検討（査読待ち初稿）}\mbox{}\par

\noindent\textbf{1：適切} 発熱反応である。式の生成側の92 kJは反応式1回分の放熱を表す。\par

\noindent\textbf{2：適切} 定温定積で窒素を加えると生成側へ進む。直後の \(Q=p_{\mathrm{NH_3}}^2/(p_{\mathrm{N_2}}p_{\mathrm{H_2}}^3)\) が小さくなる。\par

\noindent\textbf{3：不適切} 平衡後の触媒追加で生成側へ移る。触媒は平衡定数を変えず、正逆の反応を促進する。\par

\noindent\textbf{4：適切} 鉄触媒のハーバー・ボッシュ法で製造する。鉄系触媒による工業合成である。実用温度は反応速度も考慮して選ぶ。\par

\noindent\textbf{5：適切} 低温・高圧ほど生成に有利である。発熱方向は低温で平衡上有利となり、気体の化学量論数が4から2へ減る方向は圧縮で有利となる。\par

\paragraph{再利用する考え方} 反応速度と平衡位置、全圧と分圧、濃度と活量を分ける。式の適用条件を確認し、原典の不正確な記述を科学的に正しい文へ置き換えない。

\paragraph{共通理論} \hyperref[theory:chem_eq:quotient]{第\ref*{theory:chem_eq:quotient}節}, \hyperref[theory:chem_eq:ammonia]{第\ref*{theory:chem_eq:ammonia}節}, \hyperref[theory:chem_eq:catalyst]{第\ref*{theory:chem_eq:catalyst}節}

\par 参考：\cite{iupac_greenbook2012}。

\paragraph{関連問題} \hyperref[q:h30-reg-applied-16]{H30-REG-APPLIED-16}, \hyperref[q:r04-reg-applied-10]{R04-REG-APPLIED-10}, \hyperref[q:r07-reg-applied-13]{R07-REG-APPLIED-13}

\par 難易度 2（校準初稿）、優先度未評価。

\subsection{平成30年度 問16}
\label{q:h30-reg-applied-16}

\noindent\textbf{原典：} sources/H30/17\_applied\_science.pdf、PDF物理5頁。

\textbf{公式正答：} 4（single）。

\paragraph{問題の要旨} 発熱するアンモニア合成の温度・圧力・窒素追加・触媒の効果を判定する。

\Needspace{4\baselineskip}\paragraph{確認した条件・記述（独自要約）}\mbox{}\par

\noindent\textbf{確認範囲：} 問題幹、否定、全5肢、数式、数値、単位を原典画像で確認した。\par

\noindent\textbf{成立条件：} \(\mathrm{N_2+3H_2\rightleftharpoons2NH_3}+92\,\mathrm{kJ}\)。窒素追加は温度・体積一定。気体の分圧を用いる理想気体近似。\par

\paragraph{問われる核心} 発熱するアンモニア合成の温度・圧力・窒素追加・触媒の効果を判定する。

\paragraph{必要知識} \(\mathrm{N_2+3H_2\rightleftharpoons2NH_3}+92\,\mathrm{kJ}\)。窒素追加は温度・体積一定。気体の分圧を用いる理想気体近似。

\paragraph{考え方・正答を選ぶ理由（査読待ち初稿）} 反応式は窒素1、水素3、アンモニア2であり、生成側に92 kJを持つ。窒素追加は定温定積で反応商を低下させる。低温は発熱方向、高圧は気体量の少ない生成方向を平衡上有利にする。触媒による平衡移動だけが誤りであり、4を選ぶ。

\Needspace{4\baselineskip}\paragraph{選択肢の検討（査読待ち初稿）}\mbox{}\par

\noindent\textbf{1：適切} 発熱反応である。式の生成側の92 kJは反応式1回分の放熱を表す。\par

\noindent\textbf{2：適切} 定温定積で窒素を加えると生成側へ進む。直後の \(Q=p_{\mathrm{NH_3}}^2/(p_{\mathrm{N_2}}p_{\mathrm{H_2}}^3)\) が小さくなる。\par

\noindent\textbf{3：適切} 低温・高圧ほど生成に有利である。発熱方向は低温で平衡上有利となり、気体の化学量論数が4から2へ減る方向は圧縮で有利となる。\par

\noindent\textbf{4：不適切} 平衡後の触媒追加で生成側へ移る。触媒は平衡定数を変えず、正逆の反応を促進する。\par

\noindent\textbf{5：適切} 鉄触媒のハーバー・ボッシュ法で製造する。鉄系触媒による工業合成である。実用温度は反応速度も考慮して選ぶ。\par

\paragraph{再利用する考え方} 反応速度と平衡位置、全圧と分圧、濃度と活量を分ける。式の適用条件を確認し、原典の不正確な記述を科学的に正しい文へ置き換えない。

\paragraph{共通理論} \hyperref[theory:chem_eq:quotient]{第\ref*{theory:chem_eq:quotient}節}, \hyperref[theory:chem_eq:ammonia]{第\ref*{theory:chem_eq:ammonia}節}, \hyperref[theory:chem_eq:catalyst]{第\ref*{theory:chem_eq:catalyst}節}

\par 参考：\cite{iupac_greenbook2012}。

\paragraph{関連問題} \hyperref[q:h28-reg-applied-18]{H28-REG-APPLIED-18}

\par 難易度 2（校準初稿）、優先度未評価。

\subsection{令和4年度 問10}
\label{q:r04-reg-applied-10}

\noindent\textbf{原典：} sources/R04/17\_applied\_science.pdf、PDF物理3頁。

\textbf{公式正答：} 3（single）。

\paragraph{問題の要旨} 発熱するアンモニア合成の温度・圧力・窒素追加・触媒の効果を判定する。

\Needspace{4\baselineskip}\paragraph{確認した条件・記述（独自要約）}\mbox{}\par

\noindent\textbf{確認範囲：} 問題幹、否定、全5肢、数式、数値、単位を原典画像で確認した。\par

\noindent\textbf{成立条件：} \(\mathrm{N_2+3H_2\rightleftharpoons2NH_3}+92\,\mathrm{kJ}\)。窒素追加は温度・体積一定。気体の分圧を用いる理想気体近似。\par

\paragraph{問われる核心} 発熱するアンモニア合成の温度・圧力・窒素追加・触媒の効果を判定する。

\paragraph{必要知識} \(\mathrm{N_2+3H_2\rightleftharpoons2NH_3}+92\,\mathrm{kJ}\)。窒素追加は温度・体積一定。気体の分圧を用いる理想気体近似。

\paragraph{考え方・正答を選ぶ理由（査読待ち初稿）} 反応式は窒素1、水素3、アンモニア2であり、生成側に92 kJを持つ。窒素追加は定温定積で反応商を低下させる。低温は発熱方向、高圧は気体量の少ない生成方向を平衡上有利にする。触媒による平衡移動だけが誤りであり、3を選ぶ。

\Needspace{4\baselineskip}\paragraph{選択肢の検討（査読待ち初稿）}\mbox{}\par

\noindent\textbf{1：適切} 発熱反応である。式の生成側の92 kJは反応式1回分の放熱を表す。\par

\noindent\textbf{2：適切} 定温定積で窒素を加えると生成側へ進む。直後の \(Q=p_{\mathrm{NH_3}}^2/(p_{\mathrm{N_2}}p_{\mathrm{H_2}}^3)\) が小さくなる。\par

\noindent\textbf{3：不適切} 平衡後の触媒追加で生成側へ移る。触媒は平衡定数を変えず、正逆の反応を促進する。\par

\noindent\textbf{4：適切} 鉄触媒のハーバー・ボッシュ法で製造する。鉄系触媒による工業合成である。実用温度は反応速度も考慮して選ぶ。\par

\noindent\textbf{5：適切} 低温・高圧ほど生成に有利である。発熱方向は低温で平衡上有利となり、気体の化学量論数が4から2へ減る方向は圧縮で有利となる。\par

\paragraph{再利用する考え方} 反応速度と平衡位置、全圧と分圧、濃度と活量を分ける。式の適用条件を確認し、原典の不正確な記述を科学的に正しい文へ置き換えない。

\paragraph{共通理論} \hyperref[theory:chem_eq:quotient]{第\ref*{theory:chem_eq:quotient}節}, \hyperref[theory:chem_eq:ammonia]{第\ref*{theory:chem_eq:ammonia}節}, \hyperref[theory:chem_eq:catalyst]{第\ref*{theory:chem_eq:catalyst}節}

\par 参考：\cite{iupac_greenbook2012}。

\paragraph{関連問題} \hyperref[q:h28-reg-applied-18]{H28-REG-APPLIED-18}

\par 難易度 2（校準初稿）、優先度未評価。

\subsection{令和7年度 問13}
\label{q:r07-reg-applied-13}

\noindent\textbf{原典：} sources/R07/17\_applied\_science.pdf、PDF物理5頁。

\textbf{公式正答：} 3（single）。

\paragraph{問題の要旨} 四酸化二窒素の解離平衡で、反応種の分圧と不活性気体による全圧変化を区別する。

\Needspace{4\baselineskip}\paragraph{確認した条件・記述（独自要約）}\mbox{}\par

\noindent\textbf{確認範囲：} 問題幹、否定、全5肢、数式、数値、単位を原典画像で確認した。\par

\noindent\textbf{成立条件：} 1、2、3、5は温度・体積一定、4は温度一定で体積増加。理想気体近似。不活性な窒素と固体触媒は反応商に含めない。\par

\paragraph{問われる核心} 四酸化二窒素の解離平衡で、反応種の分圧と不活性気体による全圧変化を区別する。

\paragraph{必要知識} 1、2、3、5は温度・体積一定、4は温度一定で体積増加。理想気体近似。不活性な窒素と固体触媒は反応商に含めない。

\paragraph{考え方・正答を選ぶ理由（査読待ち初稿）} \(\mathrm{N_2O_4\rightleftharpoons2NO_2}\) の \(Q=p_{\mathrm{NO_2}}^2/p_{\mathrm{N_2O_4}}\) を各操作直後に比較する。一定温度・体積で窒素を加えても反応種の \(p_i=n_iRT/V\) は変わらず、Qも変わらない。全圧が上がるだけで左へ進むとする3が誤りである。体積を増やす場合は全反応種の分圧が下がり、Qは体積に反比例して下がるため右へ進む。

\Needspace{4\baselineskip}\paragraph{選択肢の検討（査読待ち初稿）}\mbox{}\par

\noindent\textbf{1：適切} 定温定積で四酸化二窒素を加えると右へ。反応商の分母が増えてQが下がる。\par

\noindent\textbf{2：適切} 定温定積で二酸化窒素を除くと右へ。反応商の分子が減ってQが下がる。\par

\noindent\textbf{3：不適切} 窒素を加えて圧力を上げると左へ。一定体積では反応種の分圧が変わらず、圧縮による圧力上昇とは異なる。\par

\noindent\textbf{4：適切} 一定温度で体積を増すと右へ。反応商が下がり気体の化学量論数の多い側へ進む。\par

\noindent\textbf{5：適切} 定温定積で固体触媒を加えても平衡は変わらない。触媒はKを変えず、反応商の因子にもならない。\par

\paragraph{再利用する考え方} 反応速度と平衡位置、全圧と分圧、濃度と活量を分ける。式の適用条件を確認し、原典の不正確な記述を科学的に正しい文へ置き換えない。

\paragraph{共通理論} \hyperref[theory:chem_eq:quotient]{第\ref*{theory:chem_eq:quotient}節}, \hyperref[theory:chem_eq:inert]{第\ref*{theory:chem_eq:inert}節}

\par 参考：\cite{iupac_greenbook2012}。

\paragraph{関連問題} \hyperref[q:h28-reg-applied-18]{H28-REG-APPLIED-18}

\par 難易度 2（校準初稿）、優先度未評価。

\section{反応ギブズエネルギー・活量・平衡定数}

\subsection{令和3年度 問10}
\label{q:r03-reg-applied-10}

\noindent\textbf{原典：} sources/R03/17\_applied\_science.pdf、PDF物理4頁。

\textbf{公式正答：} 4（single）。

\paragraph{問題の要旨} 一定温度・圧力のAとBの可逆反応で、実際の反応ギブズエネルギーと標準量を区別する。

\Needspace{4\baselineskip}\paragraph{確認した条件・記述（独自要約）}\mbox{}\par

\noindent\textbf{確認範囲：} 問題幹、否定、全5肢、数式、数値、単位を原典画像で確認した。\par

\noindent\textbf{成立条件：} 定温定圧の \(\mathrm{A\rightleftharpoons B}\)。標準生成ギブズエネルギーの差と実際の組成による反応量を区別する。濃度比で活量比を近似する条件を明記する。\par

\paragraph{問われる核心} 一定温度・圧力のAとBの可逆反応で、実際の反応ギブズエネルギーと標準量を区別する。

\paragraph{必要知識} 定温定圧の \(\mathrm{A\rightleftharpoons B}\)。標準生成ギブズエネルギーの差と実際の組成による反応量を区別する。濃度比で活量比を近似する条件を明記する。

\paragraph{考え方・正答を選ぶ理由（査読待ち初稿）} \(\Delta_rG=\Delta_rG^\circ+RT\ln Q\) であり、AとBの1対1反応では \(Q=a_B/a_A\) である。一定温度でも組成が変われば実際の反応ギブズエネルギーは変わる。平衡の \(\Delta_rG=0\) から \(K=\exp(-\Delta_rG^\circ/RT)\) が得られる。濃度比変化を無視する4が不適切である。

\Needspace{4\baselineskip}\paragraph{選択肢の検討（査読待ち初稿）}\mbox{}\par

\noindent\textbf{1：適切} 反応ギブズエネルギー負なら正方向へ進む。定温定圧で正方向の反応進行が系のギブズエネルギーを低下させる。\par

\noindent\textbf{2：適切} 平衡では反応ギブズエネルギーが0。反応進行度に関するギブズエネルギーの微分が0となる。\par

\noindent\textbf{3：適切} 標準生成量から反応標準量を求める。生成物の標準生成ギブズエネルギーの和から反応物の和を化学量論係数付きで引く。\par

\noindent\textbf{4：不適切} 一定温度なら濃度比が変わっても反応量は一定。標準反応量と異なり、実際の反応量は反応商に依存する。\par

\noindent\textbf{5：適切} 平衡定数は標準生成ギブズエネルギーから求まる。標準反応量を求めて指数式に代入する。\par

\paragraph{再利用する考え方} 反応速度と平衡位置、全圧と分圧、濃度と活量を分ける。式の適用条件を確認し、原典の不正確な記述を科学的に正しい文へ置き換えない。

\paragraph{共通理論} \hyperref[theory:chem_eq:gibbs]{第\ref*{theory:chem_eq:gibbs}節}, \hyperref[theory:chem_eq:quotient]{第\ref*{theory:chem_eq:quotient}節}

\par 参考：\cite{iupac_greenbook2012}。

\paragraph{関連問題} \hyperref[q:r05-reg-applied-17]{R05-REG-APPLIED-17}

\par 難易度 2（校準初稿）、優先度未評価。

\subsection{令和5年度 問17}
\label{q:r05-reg-applied-17}

\noindent\textbf{原典：} sources/R05/17\_applied\_science.pdf、PDF物理6頁。

\textbf{公式正答：} 1（single）。

\paragraph{問題の要旨} 一定温度・圧力のAとBの可逆反応で、実際の反応ギブズエネルギーと標準量を区別する。

\Needspace{4\baselineskip}\paragraph{確認した条件・記述（独自要約）}\mbox{}\par

\noindent\textbf{確認範囲：} 問題幹、否定、全5肢、数式、数値、単位を原典画像で確認した。\par

\noindent\textbf{成立条件：} 定温定圧の \(\mathrm{A\rightleftharpoons B}\)。標準生成ギブズエネルギーの差と実際の組成による反応量を区別する。濃度比で活量比を近似する条件を明記する。\par

\paragraph{問われる核心} 一定温度・圧力のAとBの可逆反応で、実際の反応ギブズエネルギーと標準量を区別する。

\paragraph{必要知識} 定温定圧の \(\mathrm{A\rightleftharpoons B}\)。標準生成ギブズエネルギーの差と実際の組成による反応量を区別する。濃度比で活量比を近似する条件を明記する。

\paragraph{考え方・正答を選ぶ理由（査読待ち初稿）} \(\Delta_rG=\Delta_rG^\circ+RT\ln Q\) であり、AとBの1対1反応では \(Q=a_B/a_A\) である。一定温度でも組成が変われば実際の反応ギブズエネルギーは変わる。平衡の \(\Delta_rG=0\) から \(K=\exp(-\Delta_rG^\circ/RT)\) が得られる。濃度比変化を無視する1が不適切である。

\Needspace{4\baselineskip}\paragraph{選択肢の検討（査読待ち初稿）}\mbox{}\par

\noindent\textbf{1：不適切} 一定温度なら濃度比が変わっても反応量は一定。標準反応量と異なり、実際の反応量は反応商に依存する。\par

\noindent\textbf{2：適切} 反応ギブズエネルギー負なら正方向へ進む。定温定圧で正方向の反応進行が系のギブズエネルギーを低下させる。\par

\noindent\textbf{3：適切} 平衡定数は標準生成ギブズエネルギーから求まる。標準反応量を求めて指数式に代入する。\par

\noindent\textbf{4：適切} 標準生成量から反応標準量を求める。生成物の標準生成ギブズエネルギーの和から反応物の和を化学量論係数付きで引く。\par

\noindent\textbf{5：適切} 平衡では反応ギブズエネルギーが0。反応進行度に関するギブズエネルギーの微分が0となる。\par

\paragraph{再利用する考え方} 反応速度と平衡位置、全圧と分圧、濃度と活量を分ける。式の適用条件を確認し、原典の不正確な記述を科学的に正しい文へ置き換えない。

\paragraph{共通理論} \hyperref[theory:chem_eq:gibbs]{第\ref*{theory:chem_eq:gibbs}節}, \hyperref[theory:chem_eq:quotient]{第\ref*{theory:chem_eq:quotient}節}

\par 参考：\cite{iupac_greenbook2012}。

\paragraph{関連問題} \hyperref[q:r03-reg-applied-10]{R03-REG-APPLIED-10}

\par 難易度 2（校準初稿）、優先度未評価。

\section{酸化数・電子収支}

\subsection{平成25年度 問11}
\label{q:h25-reg-applied-11}

\noindent\textbf{原典：} sources/H25/17\_applied\_science.pdf、PDF物理5頁。

\textbf{公式正答：} 1（single）。

\paragraph{問題の要旨} 反応前後の酸化数を求め、還元に当たる変化を選ぶ。

\Needspace{4\baselineskip}\paragraph{確認した条件・記述（独自要約）}\mbox{}\par

\noindent\textbf{確認範囲：} 問題幹、数値、式、否定、全5候補を原典画像で照合した。\par

\noindent\textbf{成立条件：} 単体の酸化数は0、単原子イオンは電荷に等しい。この問題の通常の酸化物でOは−2、H2SでHは+1である。式は変化の概略であって全反応の釣合い式ではない。\par

\paragraph{問われる核心} 反応前後の酸化数を求め、還元に当たる変化を選ぶ。

\paragraph{必要知識} 単体の酸化数は0、単原子イオンは電荷に等しい。この問題の通常の酸化物でOは−2、H2SでHは+1である。式は変化の概略であって全反応の釣合い式ではない。

\paragraph{考え方・正答を選ぶ理由（査読待ち初稿）} \(\mathrm{MnO_4^-}\) ではMnの酸化数が \(x+4(-2)=-1\) より+7である。\(\mathrm{Mn^{2+}}\) では+2なので5電子を受け取る還元であり1を選ぶ。他の変化は酸化数不変又は増加である。

\Needspace{4\baselineskip}\paragraph{選択肢の検討（査読待ち初稿）}\mbox{}\par

\noindent\textbf{1：適切} \(\mathrm{MnO_4^-}\to\mathrm{Mn^{2+}}\)。Mnは+7から+2へ低下し還元である。\par

\noindent\textbf{2：不適切} NaCl→NaOH。Naは両側とも+1であり、Naの還元ではない。ClはCl−として別に移るなど全反応の収支が省略された表示である。\par

\noindent\textbf{3：不適切} \(2\mathrm{I^-}\to\mathrm{I_2}\)。Iは−1から0へ増加し、2電子を放出する酸化である。\par

\noindent\textbf{4：不適切} Mg→MgCl2。Mgは0から+2へ増加し酸化である。\par

\noindent\textbf{5：不適切} H2S→S。Sは−2から0へ増加し酸化である。\par

\paragraph{再利用する考え方} 酸素の付加だけで判定せず、対象元素の酸化数と電子の受け渡しを追う。酸化数は形式的な電荷であり、共有結合内の実電荷そのものではない。

\paragraph{共通理論} \hyperref[theory:electrochem:oxidation_number]{第\ref*{theory:electrochem:oxidation_number}節}, \hyperref[theory:electrochem:half_reaction]{第\ref*{theory:electrochem:half_reaction}節}

\par 参考：\cite{mit_redox2014}。

\par 難易度 2（校準初稿）、優先度未評価。

\section{中和後の塩の加水分解とpH}

\subsection{平成24年度 問10}
\label{q:h24-reg-applied-10}

\noindent\textbf{原典：} sources/H24/17\_applied\_science.pdf、PDF物理3頁。

\textbf{公式正答：} 5（single）。

\paragraph{問題の要旨} 同濃度・同体積で中和した3組の酸塩基について、生成塩の加水分解からpHの順序を求める。

\Needspace{4\baselineskip}\paragraph{確認した条件・記述（独自要約）}\mbox{}\par

\noindent\textbf{確認範囲：} 問題幹、否定、全5肢、数式、数値、単位を原典画像で確認した。\par

\noindent\textbf{成立条件：} 各水溶液は同じモル濃度、同じ体積。1対1の中和を先に処理する。強酸強塩基塩は水の自己解離に対して中性、弱酸塩は塩基性、弱塩基塩は酸性。\par

\paragraph{問われる核心} 同濃度・同体積で中和した3組の酸塩基について、生成塩の加水分解からpHの順序を求める。

\paragraph{必要知識} 各水溶液は同じモル濃度、同じ体積。1対1の中和を先に処理する。強酸強塩基塩は水の自己解離に対して中性、弱酸塩は塩基性、弱塩基塩は酸性。

\paragraph{考え方・正答を選ぶ理由（査読待ち初稿）} Aの酢酸と水酸化ナトリウムから生じる酢酸イオンは加水分解で水酸化物イオンを生じ、塩基性となる。Bの塩酸と水酸化ナトリウムは中性の塩、Cの塩酸とアンモニアはアンモニウムイオンを含む酸性の塩である。低いpHからC、B、Aとなり5を選ぶ。

\Needspace{4\baselineskip}\paragraph{選択肢の検討（査読待ち初稿）}\mbox{}\par

\noindent\textbf{1：不適切} A・B・C。酸性C、中性B、塩基性Aの順に一致しない。\par

\noindent\textbf{2：不適切} B・C・A。酸性C、中性B、塩基性Aの順に一致しない。\par

\noindent\textbf{3：不適切} B・A・C。酸性C、中性B、塩基性Aの順に一致しない。\par

\noindent\textbf{4：不適切} C・A・B。酸性C、中性B、塩基性Aの順に一致しない。\par

\noindent\textbf{5：適切} C・B・A。酸性C、中性B、塩基性Aの順に一致する。\par

\paragraph{再利用する考え方} 反応速度と平衡位置、全圧と分圧、濃度と活量を分ける。式の適用条件を確認し、原典の不正確な記述を科学的に正しい文へ置き換えない。

\paragraph{共通理論} \hyperref[theory:chem_eq:hydrolysis]{第\ref*{theory:chem_eq:hydrolysis}節}

\par 参考：\cite{iupac_greenbook2012,mit_buffer2008}。

\paragraph{関連問題} \hyperref[q:h30-reg-applied-17]{H30-REG-APPLIED-17}

\par 難易度 2（校準初稿）、優先度未評価。

\subsection{平成30年度 問17}
\label{q:h30-reg-applied-17}

\noindent\textbf{原典：} sources/H30/17\_applied\_science.pdf、PDF物理6頁。

\textbf{公式正答：} 3（single）。

\par\textbf{検討未確定：} 形式濃度1 mol/dm3の塩溶液に活量係数1を用いる計算上の近似を明示する。原典は温度・活量係数・実在溶液の調製可能性を指定しないため、実測pHの断定とは区別する。

\paragraph{問題の要旨} 安息香酸と等モルの強塩基を混合し、生成塩の加水分解でpHを計算する。

\Needspace{4\baselineskip}\paragraph{確認した条件・記述（独自要約）}\mbox{}\par

\noindent\textbf{確認範囲：} 問題幹、否定、全5肢、数式、数値、単位を原典画像で確認した。\par

\noindent\textbf{成立条件：} \(K_a=10^{-4}\,\mathrm{mol\,dm^{-3}}\)、\(K_w=10^{-14}\,\mathrm{mol^2\,dm^{-6}}\)。体積加成、濃度で平衡式を扱う理想溶液近似、加水分解量が1 mol/dm3に比べ十分小さいこと。\par

\paragraph{問われる核心} 安息香酸と等モルの強塩基を混合し、生成塩の加水分解でpHを計算する。

\paragraph{必要知識} \(K_a=10^{-4}\,\mathrm{mol\,dm^{-3}}\)、\(K_w=10^{-14}\,\mathrm{mol^2\,dm^{-6}}\)。体積加成、濃度で平衡式を扱う理想溶液近似、加水分解量が1 mol/dm3に比べ十分小さいこと。

\paragraph{考え方・正答を選ぶ理由（査読待ち初稿）} 両溶液は \(2\,\mathrm{mol\,dm^{-3}}\)、各 \(10\,\mathrm{cm^3}=0.010\,\mathrm{dm^3}\) である。各物質量は0.020 molなので過剰の強塩基は残らない。混合体積0.020 dm3より安息香酸塩の形式濃度は1 mol/dm3である。\(K_b=K_w/K_a=10^{-10}\,\mathrm{mol\,dm^{-3}}\) とし、加水分解量 \(x\) は \(x^2/(1-x)=10^{-10}\) から約 \(10^{-5}\,\mathrm{mol\,dm^{-3}}\) となる。与えた \(K_w=10^{-14}\) に対応してpOH=\allowbreak{}5、pH=\allowbreak{}9であり3を選ぶ。

\Needspace{4\baselineskip}\paragraph{選択肢の検討（査読待ち初稿）}\mbox{}\par

\noindent\textbf{1：不適切} 7。等モル中和は過剰酸塩基を消すが、弱酸の塩の加水分解を消さない。\par

\noindent\textbf{2：不適切} 8。加水分解の平方根から得たpOH=\allowbreak{}5と対応しない。\par

\noindent\textbf{3：適切} 9。水酸化物イオン濃度が約10のマイナス5乗なので対応する。\par

\noindent\textbf{4：不適切} 10。形式濃度1とKbの積の平方根を取った結果と対応しない。\par

\noindent\textbf{5：不適切} 14。未反応の1 mol/dm3強塩基が残る場合ではない。\par

\paragraph{再利用する考え方} 反応速度と平衡位置、全圧と分圧、濃度と活量を分ける。式の適用条件を確認し、原典の不正確な記述を科学的に正しい文へ置き換えない。

\paragraph{共通理論} \hyperref[theory:chem_eq:hydrolysis]{第\ref*{theory:chem_eq:hydrolysis}節}

\par 参考：\cite{iupac_greenbook2012,mit_buffer2008}。

\paragraph{関連問題} \hyperref[q:h24-reg-applied-10]{H24-REG-APPLIED-10}

\par 難易度 3（校準初稿）、優先度未評価。

\section{溶解度積・pH・共通イオン}

\subsection{平成29年度 問13}
\label{q:h29-reg-applied-13}

\noindent\textbf{原典：} sources/H29/17\_applied\_science.pdf、PDF物理4頁。

\textbf{公式正答：} 3（single）。

\paragraph{問題の要旨} 一次可逆反応、反応性衝突、溶解度積、連鎖反応、熱暴走を比較する。

\Needspace{4\baselineskip}\paragraph{確認した条件・記述（独自要約）}\mbox{}\par

\noindent\textbf{確認範囲：} 問題幹、否定、全5肢、数式、数値、単位を原典画像で確認した。\par

\noindent\textbf{成立条件：} 1のAとBの可逆反応は双方一次。2はClとNOClの衝突。3は \(\mathrm{Mg(OH)_2}\)。4は連鎖開始・伝播・停止。5は発熱と温度上昇の正帰還。\par

\paragraph{問われる核心} 一次可逆反応、反応性衝突、溶解度積、連鎖反応、熱暴走を比較する。

\paragraph{必要知識} 1のAとBの可逆反応は双方一次。2はClとNOClの衝突。3は \(\mathrm{Mg(OH)_2}\)。4は連鎖開始・伝播・停止。5は発熱と温度上昇の正帰還。

\paragraph{考え方・正答を選ぶ理由（査読待ち初稿）} 水酸化マグネシウムの飽和溶液では \(K_{\mathrm{sp}}=a_{\mathrm{Mg^{2+}}}a_{\mathrm{OH^-}}^2\) である。一定温度でもpHにより水酸化物イオンの活量が変わるため、溶解度は一般に変化する。pH依存性を否定する3が不適切である。

\Needspace{4\baselineskip}\paragraph{選択肢の検討（査読待ち初稿）}\mbox{}\par

\noindent\textbf{1：適切} 一次正逆反応の速度定数比が平衡定数。平衡で \(k_f[A]=k_r[B]\) なので濃度基準の \(K=[B]/[A]=k_f/k_r\) となる。\par

\noindent\textbf{2：適切} 十分なエネルギーの衝突でも必ず反応するとは限らない。分子の向きや反応経路にも依存する。\par

\noindent\textbf{3：不適切} 水酸化マグネシウムの溶解度はpHに依存しない。共通イオン効果を通じて水酸化物イオンの活量が溶解量を変える。\par

\noindent\textbf{4：適切} 連鎖反応は開始・伝播・停止を含む。反応性中間体の生成、再生、消滅を区別する。\par

\noindent\textbf{5：適切} 発熱が速度を上げて熱爆発を起こし得る。発熱と除熱の釣合いが失われると温度上昇と反応加速が循環する。\par

\paragraph{再利用する考え方} 反応速度と平衡位置、全圧と分圧、濃度と活量を分ける。式の適用条件を確認し、原典の不正確な記述を科学的に正しい文へ置き換えない。

\paragraph{共通理論} \hyperref[theory:chem_eq:solubility]{第\ref*{theory:chem_eq:solubility}節}, \hyperref[theory:chem_eq:rates]{第\ref*{theory:chem_eq:rates}節}

\par 参考：\cite{iupac_greenbook2012}。

\par 難易度 2（校準初稿）、優先度未評価。

\section{飽和蒸気圧・沸騰・密封条件}

\subsection{平成24年度 問13}
\label{q:h24-reg-applied-13}

\noindent\textbf{原典：} sources/H24/17\_applied\_science.pdf、PDF物理4頁。

\textbf{公式正答：} 3（single）。

\par\textbf{検討未確定：} 原典は純物質という限定、膨張後も液体が残る条件、密封加熱時の空気・気相部を明示しない。標準的な想定での公式判定と、任意容器への全称を区別する。H24の肢2には平衡状態の明記もなく、過渡状態を排除できない。

\paragraph{問題の要旨} 飽和蒸気圧、動的平衡、定圧沸騰と密封加熱を条件ごとに判定する。

\Needspace{4\baselineskip}\paragraph{確認した条件・記述（独自要約）}\mbox{}\par

\noindent\textbf{確認範囲：} 問題幹、条件、否定、数値、単位、全5肢を原典画像で照合した。\par

\noindent\textbf{成立条件：} 定温、気液共存と平衡、定圧沸騰、室温大気圧での密封、不揮発性溶質という条件。蒸気の分圧と容器の全圧を区別する。\par

\paragraph{問われる核心} 飽和蒸気圧、動的平衡、定圧沸騰と密封加熱を条件ごとに判定する。

\paragraph{必要知識} 定温、気液共存と平衡、定圧沸騰、室温大気圧での密封、不揮発性溶質という条件。蒸気の分圧と容器の全圧を区別する。

\paragraph{考え方・正答を選ぶ理由（査読待ち初稿）} 純液体の定圧沸騰では、二相が残る間の温度は一定である。上昇し続けるとする3を選ぶ。密封容器の全圧と水の分圧を分け、定温膨張は液体が残る場合に限って飽和蒸気圧へ戻ると読む。

\Needspace{4\baselineskip}\paragraph{選択肢の検討（査読待ち初稿）}\mbox{}\par

\noindent\textbf{1：条件付きで適切} 定温で二相共存の密封容器を広げても蒸気圧は変わらない。膨張後も液体が残り平衡に達すると、追加蒸発により同じ温度の飽和蒸気圧に戻る。全液体が蒸発した後は成立しない。\par

\noindent\textbf{2：条件付きで適切} 密封容器に液体と気体が共存すると蒸発速度と凝縮速度が等しい。動的平衡に達した場合に等しい。原典はこの肢に平衡状態の語がなく、共存だけでは正味蒸発が起こる過渡状態も含まれる。\par

\noindent\textbf{3：不適切} 大気圧の沸騰中も温度は上昇し続ける。純液体の定圧気液共存では加熱が蒸発潜熱へ使われ、温度は沸点に保たれる。\par

\noindent\textbf{4：条件付きで適切} 室温・大気圧で水を入れ密封し加熱すると100°Cで沸騰しない。気相部に空気が残る定積容器なら、100°Cで水蒸気分圧が約1気圧でも全圧には空気分圧が加わる。沸騰の判定は全外圧による。空気を除いた容器などへ無条件に一般化しない。\par

\noindent\textbf{5：適切} 不揮発性溶質を溶かした希薄溶液の蒸気圧は下がる。同じ温度で溶媒の活量が低下し、理想希薄近似では \(p_A=x_Ap_A^*\) となる。溶質は蒸気圧に寄与しない。\par

\paragraph{再利用する考え方} 圧力を一定に保つ開放容器と、物質量・体積を固定する密封容器は異なる。液体が残る条件を飽和蒸気圧の議論に付ける。

\paragraph{共通理論} \hyperref[theory:phase_sep:vapor]{第\ref*{theory:phase_sep:vapor}節}, \hyperref[theory:phase_sep:boiling]{第\ref*{theory:phase_sep:boiling}節}, \hyperref[theory:phase_sep:sealed]{第\ref*{theory:phase_sep:sealed}節}, \hyperref[theory:phase_sep:raoult]{第\ref*{theory:phase_sep:raoult}節}

\par 参考：\cite{mit_solutions2008,mit_clapeyron2008}。

\paragraph{関連問題} \hyperref[q:h27-reg-applied-17]{H27-REG-APPLIED-17}

\par 難易度 2（校準初稿）、優先度未評価。

\subsection{平成27年度 問17}
\label{q:h27-reg-applied-17}

\noindent\textbf{原典：} sources/H27/17\_applied\_science.pdf、PDF物理6頁。

\textbf{公式正答：} 3（single）。

\par\textbf{検討未確定：} 原典は純物質という限定、膨張後も液体が残る条件、密封加熱時の空気・気相部を明示しない。標準的な想定での公式判定と、任意容器への全称を区別する。

\paragraph{問題の要旨} 飽和蒸気圧、動的平衡、定圧沸騰と密封加熱を条件ごとに判定する。

\Needspace{4\baselineskip}\paragraph{確認した条件・記述（独自要約）}\mbox{}\par

\noindent\textbf{確認範囲：} 問題幹、条件、否定、数値、単位、全5肢を原典画像で照合した。\par

\noindent\textbf{成立条件：} 定温、気液共存と平衡、定圧沸騰、室温大気圧での密封、不揮発性溶質という条件。蒸気の分圧と容器の全圧を区別する。\par

\paragraph{問われる核心} 飽和蒸気圧、動的平衡、定圧沸騰と密封加熱を条件ごとに判定する。

\paragraph{必要知識} 定温、気液共存と平衡、定圧沸騰、室温大気圧での密封、不揮発性溶質という条件。蒸気の分圧と容器の全圧を区別する。

\paragraph{考え方・正答を選ぶ理由（査読待ち初稿）} 純液体の定圧沸騰では、二相が残る間の温度は一定である。上昇し続けるとする3を選ぶ。密封容器の全圧と水の分圧を分け、定温膨張は液体が残る場合に限って飽和蒸気圧へ戻ると読む。

\Needspace{4\baselineskip}\paragraph{選択肢の検討（査読待ち初稿）}\mbox{}\par

\noindent\textbf{1：条件付きで適切} 定温で二相共存の密封容器を広げても蒸気圧は変わらない。膨張後も液体が残り平衡に達すると、追加蒸発により同じ温度の飽和蒸気圧に戻る。全液体が蒸発した後は成立しない。\par

\noindent\textbf{2：適切} 平衡の液体と気体では蒸発速度と凝縮速度が等しい。界面を往復する移動が続いても、物質量の正味変化は0となる。平衡を移動停止と混同しない。\par

\noindent\textbf{3：不適切} 大気圧の沸騰中も温度は上昇し続ける。純液体の定圧気液共存では加熱が蒸発潜熱へ使われ、温度は沸点に保たれる。\par

\noindent\textbf{4：条件付きで適切} 室温・大気圧で水を入れ密封し加熱すると100°Cで沸騰しない。気相部に空気が残る定積容器なら、100°Cで水蒸気分圧が約1気圧でも全圧には空気分圧が加わる。沸騰の判定は全外圧による。空気を除いた容器などへ無条件に一般化しない。\par

\noindent\textbf{5：適切} 不揮発性溶質を溶かした希薄溶液の蒸気圧は下がる。同じ温度で溶媒の活量が低下し、理想希薄近似では \(p_A=x_Ap_A^*\) となる。溶質は蒸気圧に寄与しない。\par

\paragraph{再利用する考え方} 圧力を一定に保つ開放容器と、物質量・体積を固定する密封容器は異なる。液体が残る条件を飽和蒸気圧の議論に付ける。

\paragraph{共通理論} \hyperref[theory:phase_sep:vapor]{第\ref*{theory:phase_sep:vapor}節}, \hyperref[theory:phase_sep:boiling]{第\ref*{theory:phase_sep:boiling}節}, \hyperref[theory:phase_sep:sealed]{第\ref*{theory:phase_sep:sealed}節}, \hyperref[theory:phase_sep:raoult]{第\ref*{theory:phase_sep:raoult}節}

\par 参考：\cite{mit_solutions2008,mit_clapeyron2008}。

\paragraph{関連問題} \hyperref[q:h24-reg-applied-13]{H24-REG-APPLIED-13}, \hyperref[q:h29-reg-applied-14]{H29-REG-APPLIED-14}, \hyperref[q:r01-reg-applied-10]{R01-REG-APPLIED-10}

\par 難易度 2（校準初稿）、優先度未評価。

\subsection{平成29年度 問14}
\label{q:h29-reg-applied-14}

\noindent\textbf{原典：} sources/H29/17\_applied\_science.pdf、PDF物理5頁。

\textbf{公式正答：} 3（single）。

\par\textbf{検討未確定：} 原典は純物質という限定、膨張後も液体が残る条件、密封加熱時の空気・気相部を明示しない。標準的な想定での公式判定と、任意容器への全称を区別する。

\paragraph{問題の要旨} 飽和蒸気圧、動的平衡、定圧沸騰と密封加熱を条件ごとに判定する。

\Needspace{4\baselineskip}\paragraph{確認した条件・記述（独自要約）}\mbox{}\par

\noindent\textbf{確認範囲：} 問題幹、条件、否定、数値、単位、全5肢を原典画像で照合した。\par

\noindent\textbf{成立条件：} 定温、気液共存と平衡、定圧沸騰、室温大気圧での密封、不揮発性溶質という条件。蒸気の分圧と容器の全圧を区別する。\par

\paragraph{問われる核心} 飽和蒸気圧、動的平衡、定圧沸騰と密封加熱を条件ごとに判定する。

\paragraph{必要知識} 定温、気液共存と平衡、定圧沸騰、室温大気圧での密封、不揮発性溶質という条件。蒸気の分圧と容器の全圧を区別する。

\paragraph{考え方・正答を選ぶ理由（査読待ち初稿）} 純液体の定圧沸騰では、二相が残る間の温度は一定である。上昇し続けるとする3を選ぶ。密封容器の全圧と水の分圧を分け、定温膨張は液体が残る場合に限って飽和蒸気圧へ戻ると読む。

\Needspace{4\baselineskip}\paragraph{選択肢の検討（査読待ち初稿）}\mbox{}\par

\noindent\textbf{1：条件付きで適切} 定温で二相共存の密封容器を広げても蒸気圧は変わらない。膨張後も液体が残り平衡に達すると、追加蒸発により同じ温度の飽和蒸気圧に戻る。全液体が蒸発した後は成立しない。\par

\noindent\textbf{2：適切} 平衡の液体と気体では蒸発速度と凝縮速度が等しい。界面を往復する移動が続いても、物質量の正味変化は0となる。平衡を移動停止と混同しない。\par

\noindent\textbf{3：不適切} 大気圧の沸騰中も温度は上昇し続ける。純液体の定圧気液共存では加熱が蒸発潜熱へ使われ、温度は沸点に保たれる。\par

\noindent\textbf{4：条件付きで適切} 室温・大気圧で水を入れ密封し加熱すると100°Cで沸騰しない。気相部に空気が残る定積容器なら、100°Cで水蒸気分圧が約1気圧でも全圧には空気分圧が加わる。沸騰の判定は全外圧による。空気を除いた容器などへ無条件に一般化しない。\par

\noindent\textbf{5：適切} 不揮発性溶質を溶かした希薄溶液の蒸気圧は下がる。同じ温度で溶媒の活量が低下し、理想希薄近似では \(p_A=x_Ap_A^*\) となる。溶質は蒸気圧に寄与しない。\par

\paragraph{再利用する考え方} 圧力を一定に保つ開放容器と、物質量・体積を固定する密封容器は異なる。液体が残る条件を飽和蒸気圧の議論に付ける。

\paragraph{共通理論} \hyperref[theory:phase_sep:vapor]{第\ref*{theory:phase_sep:vapor}節}, \hyperref[theory:phase_sep:boiling]{第\ref*{theory:phase_sep:boiling}節}, \hyperref[theory:phase_sep:sealed]{第\ref*{theory:phase_sep:sealed}節}, \hyperref[theory:phase_sep:raoult]{第\ref*{theory:phase_sep:raoult}節}

\par 参考：\cite{mit_solutions2008,mit_clapeyron2008}。

\paragraph{関連問題} \hyperref[q:h27-reg-applied-17]{H27-REG-APPLIED-17}

\par 難易度 2（校準初稿）、優先度未評価。

\subsection{令和元年度 問10}
\label{q:r01-reg-applied-10}

\noindent\textbf{原典：} sources/R01/17\_applied\_science.pdf、PDF物理3頁。

\textbf{公式正答：} 4（single）。

\par\textbf{検討未確定：} 原典は純物質という限定、膨張後も液体が残る条件、密封加熱時の空気・気相部を明示しない。標準的な想定での公式判定と、任意容器への全称を区別する。

\paragraph{問題の要旨} 飽和蒸気圧、動的平衡、定圧沸騰と密封加熱を条件ごとに判定する。

\Needspace{4\baselineskip}\paragraph{確認した条件・記述（独自要約）}\mbox{}\par

\noindent\textbf{確認範囲：} 問題幹、条件、否定、数値、単位、全5肢を原典画像で照合した。\par

\noindent\textbf{成立条件：} 定温、気液共存と平衡、定圧沸騰、室温大気圧での密封、不揮発性溶質という条件。蒸気の分圧と容器の全圧を区別する。\par

\paragraph{問われる核心} 飽和蒸気圧、動的平衡、定圧沸騰と密封加熱を条件ごとに判定する。

\paragraph{必要知識} 定温、気液共存と平衡、定圧沸騰、室温大気圧での密封、不揮発性溶質という条件。蒸気の分圧と容器の全圧を区別する。

\paragraph{考え方・正答を選ぶ理由（査読待ち初稿）} 二相共存が続く純物質の飽和蒸気圧は温度で決まる。体積を増すと液体が蒸発して気相の量が増え、平衡圧力を維持する。4を不適切とする公式判定はこの条件下で成立する。液体が尽きるほどの膨張とは区別する。

\Needspace{4\baselineskip}\paragraph{選択肢の検討（査読待ち初稿）}\mbox{}\par

\noindent\textbf{1：条件付きで適切} 室温・大気圧で水を入れ密封し加熱すると100°Cで沸騰しない。気相部に空気が残る定積容器なら、100°Cで水蒸気分圧が約1気圧でも全圧には空気分圧が加わる。沸騰の判定は全外圧による。空気を除いた容器などへ無条件に一般化しない。\par

\noindent\textbf{2：適切} 不揮発性溶質を溶かした希薄溶液の蒸気圧は下がる。同じ温度で溶媒の活量が低下し、理想希薄近似では \(p_A=x_Ap_A^*\) となる。溶質は蒸気圧に寄与しない。\par

\noindent\textbf{3：条件付きで適切} 大気圧で沸騰中の液体の温度は上昇しない。純物質の定圧で気液が共存する間は、加熱が潜熱となる。多成分混合物の組成変化や全液体の蒸発後まで一定とはしない。\par

\noindent\textbf{4：不適切} 定温の二相共存容器を広げると蒸気圧が変化する。膨張後も二相が残る標準的な想定では、同じ温度の飽和蒸気圧に戻るため変化しない。\par

\noindent\textbf{5：適切} 平衡の液体と気体では蒸発速度と凝縮速度が等しい。界面を往復する移動が続いても、物質量の正味変化は0となる。平衡を移動停止と混同しない。\par

\paragraph{再利用する考え方} 圧力を一定に保つ開放容器と、物質量・体積を固定する密封容器は異なる。液体が残る条件を飽和蒸気圧の議論に付ける。

\paragraph{共通理論} \hyperref[theory:phase_sep:vapor]{第\ref*{theory:phase_sep:vapor}節}, \hyperref[theory:phase_sep:boiling]{第\ref*{theory:phase_sep:boiling}節}, \hyperref[theory:phase_sep:sealed]{第\ref*{theory:phase_sep:sealed}節}, \hyperref[theory:phase_sep:raoult]{第\ref*{theory:phase_sep:raoult}節}

\par 参考：\cite{mit_solutions2008,mit_clapeyron2008}。

\paragraph{関連問題} \hyperref[q:h27-reg-applied-17]{H27-REG-APPLIED-17}, \hyperref[q:r04-reg-applied-14]{R04-REG-APPLIED-14}

\par 難易度 2（校準初稿）、優先度未評価。

\subsection{令和3年度 問11}
\label{q:r03-reg-applied-11}

\noindent\textbf{原典：} sources/R03/17\_applied\_science.pdf、PDF物理4頁。

\textbf{公式正答：} 5（single）。

\paragraph{問題の要旨} 純物質の蒸発と沸騰を区別し、平衡と相変化熱の関係を判定する。

\Needspace{4\baselineskip}\paragraph{確認した条件・記述（独自要約）}\mbox{}\par

\noindent\textbf{確認範囲：} 問題幹、条件、否定、数値、単位、全5肢を原典画像で照合した。\par

\noindent\textbf{成立条件：} 純物質の気液相変化。平衡で \(\Delta S_{\mathrm{vap}}=\Delta H_{\mathrm{vap}}/T\)、同じ端点間の逆変化でエンタルピー差の符号が逆になる。クラウジウス＝クラペイロン式は蒸気を理想気体とし液体体積を無視する近似。\par

\paragraph{問われる核心} 純物質の蒸発と沸騰を区別し、平衡と相変化熱の関係を判定する。

\paragraph{必要知識} 純物質の気液相変化。平衡で \(\Delta S_{\mathrm{vap}}=\Delta H_{\mathrm{vap}}/T\)、同じ端点間の逆変化でエンタルピー差の符号が逆になる。クラウジウス＝クラペイロン式は蒸気を理想気体とし液体体積を無視する近似。

\paragraph{考え方・正答を選ぶ理由（査読待ち初稿）} 沸騰は液体内部にも蒸気泡が成長できる条件であるが、表面からの蒸発は沸点未満でも起こる。液体が安定な条件にも飽和蒸気が共存し得る。沸点未満では気化しないとする5を選ぶ。

\Needspace{4\baselineskip}\paragraph{選択肢の検討（査読待ち初稿）}\mbox{}\par

\noindent\textbf{1：条件付きで適切} 蒸気圧の温度依存性はクラウジウス＝クラペイロン式で表す。液相体積を無視し蒸気を理想気体とする近似である。より一般的な関係はクラペイロン式であり、積分時は潜熱一定などの追加仮定が必要となる。\par

\noindent\textbf{2：適切} 蒸発エントロピーは蒸発エンタルピーと気化温度で表せる。平衡での可逆相変化に対し、エントロピー変化は潜熱を絶対温度で割った量である。\par

\noindent\textbf{3：適切} 平衡気液の化学ポテンシャルは等しい。相間の物質移動に対する自由エネルギー変化が0である。\par

\noindent\textbf{4：適切} 蒸発・凝縮エンタルピーは等大逆符号。同じ温度・圧力の端点を逆向きに結ぶ場合の状態関数差である。\par

\noindent\textbf{5：不適切} 沸点に満たない液体は気化しない。蒸発は表面から起こり、沸騰の開始条件を必要としない。\par

\paragraph{再利用する考え方} 蒸発、沸騰、凝縮を場所と圧力条件で区別する。相の安定性が液体優位でも、気相の物質量が必ず0になるわけではない。

\paragraph{共通理論} \hyperref[theory:phase_sep:chemical_potential]{第\ref*{theory:phase_sep:chemical_potential}節}, \hyperref[theory:phase_sep:boiling]{第\ref*{theory:phase_sep:boiling}節}, \hyperref[theory:phase_sep:latent]{第\ref*{theory:phase_sep:latent}節}

\par 参考：\cite{mit_clapeyron2008}。

\par 難易度 2（校準初稿）、優先度未評価。

\subsection{令和4年度 問14}
\label{q:r04-reg-applied-14}

\noindent\textbf{原典：} sources/R04/17\_applied\_science.pdf、PDF物理5頁。

\textbf{公式正答：} 2（single）。

\par\textbf{検討未確定：} 原典は純物質という限定、膨張後も液体が残る条件、密封加熱時の空気・気相部を明示しない。標準的な想定での公式判定と、任意容器への全称を区別する。

\paragraph{問題の要旨} 飽和蒸気圧、動的平衡、定圧沸騰と密封加熱を条件ごとに判定する。

\Needspace{4\baselineskip}\paragraph{確認した条件・記述（独自要約）}\mbox{}\par

\noindent\textbf{確認範囲：} 問題幹、条件、否定、数値、単位、全5肢を原典画像で照合した。\par

\noindent\textbf{成立条件：} 定温、気液共存と平衡、定圧沸騰、室温大気圧での密封、不揮発性溶質という条件。蒸気の分圧と容器の全圧を区別する。\par

\paragraph{問われる核心} 飽和蒸気圧、動的平衡、定圧沸騰と密封加熱を条件ごとに判定する。

\paragraph{必要知識} 定温、気液共存と平衡、定圧沸騰、室温大気圧での密封、不揮発性溶質という条件。蒸気の分圧と容器の全圧を区別する。

\paragraph{考え方・正答を選ぶ理由（査読待ち初稿）} 体積2倍と十分な静置だけでは液体残存の指定にならないが、二相共存が続く通常の読みにおいて飽和蒸気圧は変わらない。2が公式の不適切肢である。3は開放容器と明示され、密封加熱時の全圧上昇と区別できる。

\Needspace{4\baselineskip}\paragraph{選択肢の検討（査読待ち初稿）}\mbox{}\par

\noindent\textbf{1：条件付きで適切} 室温・大気圧で水を入れ密封し加熱すると100°Cで沸騰しない。気相部に空気が残る定積容器なら、100°Cで水蒸気分圧が約1気圧でも全圧には空気分圧が加わる。沸騰の判定は全外圧による。空気を除いた容器などへ無条件に一般化しない。\par

\noindent\textbf{2：不適切} 定温で容器体積を2倍にして十分静置すると蒸気圧が変わる。2倍後も液体が残れば追加蒸発で元の飽和蒸気圧へ戻る。静置時間は平衡到達を保証しても液体の残存量を保証しない。\par

\noindent\textbf{3：条件付きで適切} 大気圧で沸騰中の液体の温度は上昇しない。純物質の定圧で気液が共存する間は、加熱が潜熱となる。多成分混合物の組成変化や全液体の蒸発後まで一定とはしない。\par

\noindent\textbf{4：適切} 平衡の液体と気体では蒸発速度と凝縮速度が等しい。界面を往復する移動が続いても、物質量の正味変化は0となる。平衡を移動停止と混同しない。\par

\noindent\textbf{5：適切} 不揮発性溶質を溶かした希薄溶液の蒸気圧は下がる。同じ温度で溶媒の活量が低下し、理想希薄近似では \(p_A=x_Ap_A^*\) となる。溶質は蒸気圧に寄与しない。\par

\paragraph{再利用する考え方} 圧力を一定に保つ開放容器と、物質量・体積を固定する密封容器は異なる。液体が残る条件を飽和蒸気圧の議論に付ける。

\paragraph{共通理論} \hyperref[theory:phase_sep:vapor]{第\ref*{theory:phase_sep:vapor}節}, \hyperref[theory:phase_sep:boiling]{第\ref*{theory:phase_sep:boiling}節}, \hyperref[theory:phase_sep:sealed]{第\ref*{theory:phase_sep:sealed}節}, \hyperref[theory:phase_sep:raoult]{第\ref*{theory:phase_sep:raoult}節}

\par 参考：\cite{mit_solutions2008,mit_clapeyron2008}。

\paragraph{関連問題} \hyperref[q:r01-reg-applied-10]{R01-REG-APPLIED-10}

\par 難易度 2（校準初稿）、優先度未評価。

\section{水の相変化・過冷却}

\subsection{令和2年度 問16}
\label{q:r02-reg-applied-16}

\noindent\textbf{原典：} sources/R02/17\_applied\_science.pdf、PDF物理6頁。

\textbf{公式正答：} 5（single）。

\par\textbf{検討未確定：} 比誘電率の温度・周波数条件が省略され、三重点の0.01°Cがどの温度尺度か明記されない。過冷却の反例と公式正答は一致するが、全肢の限定を保持する。

\paragraph{問題の要旨} 水の誘電性・三重点・臨界状態と、過冷却という準安定状態を区別する。

\Needspace{4\baselineskip}\paragraph{確認した条件・記述（独自要約）}\mbox{}\par

\noindent\textbf{確認範囲：} 問題幹、条件、否定、数値、単位、全5肢を原典画像で照合した。\par

\noindent\textbf{成立条件：} 純水、大気圧、臨界圧力と温度の両指定。誘電率比較は通常の室温静的値、三重点はITS-90の慣用値として読む。\par

\paragraph{問われる核心} 水の誘電性・三重点・臨界状態と、過冷却という準安定状態を区別する。

\paragraph{必要知識} 純水、大気圧、臨界圧力と温度の両指定。誘電率比較は通常の室温静的値、三重点はITS-90の慣用値として読む。

\paragraph{考え方・正答を選ぶ理由（査読待ち初稿）} 大気圧で0°C未満の水が直ちに全て氷になるとは限らない。核生成には障壁があり、準安定の液体が過冷却として残る。したがって過冷却を否定する5が不適切である。三重点の常用温度とSIでの温度の定義は区別する。

\Needspace{4\baselineskip}\paragraph{選択肢の検討（査読待ち初稿）}\mbox{}\par

\noindent\textbf{1：条件付きで適切} 純水の比誘電率は多くの有機溶媒より大きい。室温付近の静的な比誘電率の典型比較として読む。全有機溶媒・全温度への順位ではない。\par

\noindent\textbf{2：適切} 固体が液相を経ず気体に昇華することがある。固相と気相が直接関わる相変化であり、氷も温度・圧力条件に応じて昇華する。\par

\noindent\textbf{3：適切} 純水が臨界圧力・臨界温度になると気液密度が一致し界面が消える。指定された臨界状態では共存する二相の差が消える。任意の定積容器を加熱すれば必ずその状態になるという命題ではない。\par

\noindent\textbf{4：条件付きで適切} 純水の三重点は0.01°Cである。ITS-90における慣用温度として正しい。2019年以降のSIでは水の三重点を厳密な熱力学温度の定義に用いない。\par

\noindent\textbf{5：不適切} 大気圧で液体の水を氷点下まで冷やせない。氷の核生成が遅れれば、平衡では固体が安定な温度でも準安定な液体が残り、過冷却となる。\par

\paragraph{再利用する考え方} 平衡相図は最も安定な相を示すが、相転移が起こるまでの時間を指定しない。温度尺度の慣用値と、熱力学温度の定義を分離する。

\paragraph{共通理論} \hyperref[theory:phase_sep:water]{第\ref*{theory:phase_sep:water}節}, \hyperref[theory:phase_sep:critical]{第\ref*{theory:phase_sep:critical}節}, \hyperref[theory:phase_sep:metastable]{第\ref*{theory:phase_sep:metastable}節}

\par 参考：\cite{iapws_supercooled2015,iapws95_2018,si_brochure9}。

\paragraph{関連問題} \hyperref[q:r06-reg-applied-10]{R06-REG-APPLIED-10}

\par 難易度 2（校準初稿）、優先度未評価。

\subsection{令和6年度 問10}
\label{q:r06-reg-applied-10}

\noindent\textbf{原典：} sources/R06/17\_applied\_science.pdf、PDF物理3頁。

\textbf{公式正答：} 1（single）。

\par\textbf{検討未確定：} 比誘電率の温度・周波数条件が省略され、三重点の0.01°Cがどの温度尺度か明記されない。過冷却の反例と公式正答は一致するが、全肢の限定を保持する。

\paragraph{問題の要旨} 水の誘電性・三重点・臨界状態と、過冷却という準安定状態を区別する。

\Needspace{4\baselineskip}\paragraph{確認した条件・記述（独自要約）}\mbox{}\par

\noindent\textbf{確認範囲：} 問題幹、条件、否定、数値、単位、全5肢を原典画像で照合した。\par

\noindent\textbf{成立条件：} 純水、大気圧、臨界圧力と温度の両指定。誘電率比較は通常の室温静的値、三重点はITS-90の慣用値として読む。\par

\paragraph{問われる核心} 水の誘電性・三重点・臨界状態と、過冷却という準安定状態を区別する。

\paragraph{必要知識} 純水、大気圧、臨界圧力と温度の両指定。誘電率比較は通常の室温静的値、三重点はITS-90の慣用値として読む。

\paragraph{考え方・正答を選ぶ理由（査読待ち初稿）} 大気圧で0°C未満の水が直ちに全て氷になるとは限らない。核生成には障壁があり、準安定の液体が過冷却として残る。したがって過冷却を否定する1が不適切である。三重点の常用温度とSIでの温度の定義は区別する。

\Needspace{4\baselineskip}\paragraph{選択肢の検討（査読待ち初稿）}\mbox{}\par

\noindent\textbf{1：不適切} 大気圧で液体の水を氷点下まで冷やせない。氷の核生成が遅れれば、平衡では固体が安定な温度でも準安定な液体が残り、過冷却となる。\par

\noindent\textbf{2：条件付きで適切} 純水の三重点は0.01°Cである。ITS-90における慣用温度として正しい。2019年以降のSIでは水の三重点を厳密な熱力学温度の定義に用いない。\par

\noindent\textbf{3：適切} 固体が液相を経ず気体に昇華することがある。固相と気相が直接関わる相変化であり、氷も温度・圧力条件に応じて昇華する。\par

\noindent\textbf{4：条件付きで適切} 純水の比誘電率は多くの有機溶媒より大きい。室温付近の静的な比誘電率の典型比較として読む。全有機溶媒・全温度への順位ではない。\par

\noindent\textbf{5：適切} 純水が臨界圧力・臨界温度になると気液密度が一致し界面が消える。指定された臨界状態では共存する二相の差が消える。任意の定積容器を加熱すれば必ずその状態になるという命題ではない。\par

\paragraph{再利用する考え方} 平衡相図は最も安定な相を示すが、相転移が起こるまでの時間を指定しない。温度尺度の慣用値と、熱力学温度の定義を分離する。

\paragraph{共通理論} \hyperref[theory:phase_sep:water]{第\ref*{theory:phase_sep:water}節}, \hyperref[theory:phase_sep:critical]{第\ref*{theory:phase_sep:critical}節}, \hyperref[theory:phase_sep:metastable]{第\ref*{theory:phase_sep:metastable}節}

\par 参考：\cite{iapws_supercooled2015,iapws95_2018,si_brochure9}。

\paragraph{関連問題} \hyperref[q:r02-reg-applied-16]{R02-REG-APPLIED-16}

\par 難易度 2（校準初稿）、優先度未評価。
